\documentclass[11pt]{article}

\usepackage[margin=1in]{geometry}
\usepackage{amsmath}
\usepackage{amssymb}
\usepackage{array}
\usepackage{booktabs}
\usepackage{enumitem}
\usepackage{graphicx}
\usepackage{hyperref}
\usepackage[numbers,sort&compress]{natbib}

\hypersetup{
  colorlinks=true,
  linkcolor=blue,
  citecolor=blue,
  urlcolor=blue
}

\title{Estimating SME Banking Wallet Size and Product Propensity:\\
A Survey of Relevant Literature and Modeling Approaches}
\author{Chak Wong}
\date{\today}

\begin{document}
\maketitle

\begin{abstract}
Retail banks that serve small and medium-sized enterprises (SMEs) often observe
only a partial view of a business customer's financial life. A bank may know its
own deposits, card spend, loans, treasury activity, and relationship history, but
it rarely observes the customer's balances and product holdings at competing
institutions. This creates a practical ``wallet size'' problem: estimate the
products and balances an SME holds with competitors, with deposits as the main
business object, and estimate how much of that wallet can realistically be moved
to the focal bank. This survey
reviews literatures that bear on the problem: cross-selling and next-product-to-buy
models in financial services, share-of-wallet and customer-share models,
structural bank deposit demand and bank competition, SME relationship lending,
and causal response/uplift modeling. The main conclusion is that the exact
end-to-end problem is not a single canonical academic model. It is better viewed
as a composite decision system that combines latent wallet reconstruction,
deposit and product demand estimation, offer-response modeling, and constrained
profit optimization.
\end{abstract}

\section{Introduction}

For a US retail bank serving SMEs, the commercially important question is
deposit-led wallet opportunity. The bank wants to estimate the total financial
wallet of the business, how much of that wallet sits with competitors, and how
much of the competitor-held wallet is realistically movable to the focal bank.
Deposits are the anchor product because they provide funding,
transactional information, and a relationship platform for payments, lending,
and treasury services. Other products matter, but the main use case in this
paper is deposit-led wallet sizing rather than product-assortment design.

The literature does not usually call this problem ``wallet size estimation.''
Relevant work appears under several labels: cross-selling, cross-buying,
next-product-to-buy, customer share, share-of-wallet, bank deposit demand,
bank competition, and relationship lending. The practical modeling problem is
therefore interdisciplinary. Marketing science contributes models of customer
product ownership and campaign response. Industrial organization and banking
economics contribute models of bank choice, deposit demand, and competition.
SME finance contributes theories of relationship lending and the informational
role of transaction accounts. Causal inference contributes methods for separating
who is likely to buy from who is likely to buy because of a marketing action.

This paper surveys these related literatures, organizes the model architectures
that they imply, and then recommends an integrated architecture for a bank that
wants to estimate competitor deposit wallet size and capture potential.

\section{The SME Wallet-Size Problem}

Let $i$ index an SME firm, $b$ index a bank, and $p$ index products such as
demand deposits, savings deposits, credit cards, commercial real estate loans,
term loans, lines of credit, merchant services, payroll, treasury management,
and investment funds.
Let
\[
  y_{ibp} \in \{0,1\}
\]
denote whether firm $i$ holds product $p$ at bank $b$, and let
\[
  a_{ibp} \geq 0
\]
denote the associated size, such as average deposit balance, outstanding loan
balance, card spend, or assets under management. The total product-specific
wallet is
\[
  W_{ip} = \sum_b a_{ibp}.
\]
For the focal bank $f$, the observed part is $a_{ifp}$ and $y_{ifp}$. The
competitor component
\[
  W_{ip}^{-f} = \sum_{b \neq f} a_{ibp}
\]
is generally missing or only indirectly observed.

In deployment, these objects should carry a time index because both product
need and product size move over time. Write
\[
  A_{ipt}^{-f}=\sum_{b\neq f}a_{ibpt}
\]
for the competitor-held amount of product $p$ at decision time $t$. The
business question is not only whether $A_{ipt}^{-f}$ is large. It is how much of
that amount is addressable by a specific offer. Define
\[
  A_{ipo}^{\mathrm{addr}}(H)
  =
  M_{ipo}(H)A_{ipt}^{-f},
  \qquad 0\leq M_{ipo}(H)\leq 1,
\]
where $M_{ipo}(H)$ is the structural offer- and horizon-specific movable
fraction. This is not a generic purchase propensity. It is the fraction of
competitor-held balances, borrowing, spend, or assets that can plausibly
migrate to the focal bank under offer $o$ over horizon $H$. Operational reports
may also show ratios such as moved dollars divided by estimated
competitor-held dollars; those reporting ratios should be distinguished from
the latent structural fraction $M_{ipo}(H)$.

The decision problem has three parts. First, the bank wants an estimate of
latent wallet state:
\[
  \Pr(y_{ibp}=1 \mid X_i,Z_{bt},A_{bpt},M_t), \qquad
  \mathbb{E}[a_{ibp} \mid X_i,Z_{bt},A_{bpt},M_t,y_{ibp}=1],
\]
especially for competitor banks $b \neq f$, where $Z_{bt}$ and $A_{bpt}$
summarize bank and bank-product attributes observed or imputed by the model.
Second, it wants the distribution of product amounts already held at
competitors, $A_{ipt}^{-f}$, and the addressable amount
$A_{ipo}^{\mathrm{addr}}(H)$. Third, for an offer or marketing action
$o$---for example a deposit-rate offer, relationship-manager outreach,
onboarding support, or operating-account proposal---the bank wants
outcome-specific incremental effects. For a generic outcome $Y$,
\[
  \tau_{ipo}^{Y}
  =
  \mathbb{E}[Y_{ip}(o)-Y_{ip}(0)\mid \mathcal{I}_{it}],
\]
where $Y$ may be product ownership, balance dollars, spend, risk-adjusted
relationship value, runoff, or retention. Balance lift, moved competitor
dollars, and value lift are therefore related but distinct estimands.

The distinction between prediction and causal response is crucial. A model can
predict that a firm is likely to buy a loan or maintain large deposits without
showing that a campaign will cause the firm to move wallet to the focal bank.
It can also estimate a large competitor-held wallet without identifying how much
of that wallet is movable. The literature reviewed below supports different
parts of this system.

\section{Literature Survey}

This section groups the main bodies of literature relevant to SME wallet
estimation. The key point is that the literature does not support a single
static ``wallet size'' object. It supports several linked objects: latent
product need, observed product ownership, customer share, bank choice, dynamic
liquidity demand, and causal marketing response.

\subsection{Cross-Selling and Product-Ownership Models}

The closest direct literature is financial-services cross-selling. Early work
by \citet{kamakura1991latent} applies latent trait analysis to evaluate
prospects for cross-selling financial services. This line of work treats
observed product holdings as noisy signals of latent customer needs, financial
maturity, or product propensity. For a bank, this is directly relevant because
competitor holdings are often latent and must be inferred from partial signals.

\citet{kamakura2003database} extend the database-marketing perspective using a
mixed-data factor-analyzer approach for data augmentation and prediction. The
important practical insight is that product holdings, demographics, and monetary
variables can be modeled jointly, so that missing or unobserved product
positions are not handled as independent binary classification problems. This
matters for SME wallet estimation because deposits, lending, merchant services,
and cards are correlated through firm size, cash-flow volatility, owner wealth,
industry, life-cycle stage, and bank relationship depth.

Another directly relevant stream is next-product-to-buy (NPTB) modeling.
\citet{knott2002nptb} describe NPTB models for cross-selling applications and
evaluate their use in a field setting. In banking, an NPTB model can rank likely
products for each SME relationship. The limitation is that a next-product model
usually optimizes product targeting, not total competitor wallet reconstruction.
It answers ``what is the next product to offer?'' more directly than ``how large
is the competitor-held deposit wallet?''

\citet{li2005sequential} study cross-selling sequentially ordered products in
consumer banking services using a structural multivariate probit framework. The
paper is especially useful because banking products may follow life-cycle or
business-maturity sequences: transaction accounts precede cards and merchant
services; operating deposits may precede credit lines; treasury services may
follow scale and complexity. For SME banking, the sequence is not identical to
household financial maturity, but the modeling principle carries over: product
ownership is dynamic and correlated across products.

\subsection{Share-of-Wallet and Customer-Share Models}

The share-of-wallet literature is relevant when the focal bank already has a
relationship with the SME and wants to estimate the fraction of total financial
activity captured by the bank. \citet{verhoef2003crm} examines how customer
relationship perceptions and marketing instruments affect retention and
customer share development over time. This literature separates retention from
depth: keeping a customer is not the same as gaining a larger share of that
customer's wallet.

\citet{keiningham2003sow} study customer satisfaction and actual share-of-wallet
in a business-to-business environment. This is relevant because SME banking is
closer to relationship-managed B2B service than to anonymous retail product
purchase. Satisfaction, relationship quality, and switching costs can affect
whether an SME consolidates deposits, lending, cards, and treasury activity at
one bank.

For the bank's use case, customer-share models can be written as
\[
  s_{ipf} = \frac{a_{ifp}}{W_{ip}},
\]
where $s_{ipf}$ is the focal bank's share of the SME's product-$p$ wallet.
When $W_{ip}$ is unknown, the model must combine internal balances with external
signals, market benchmarks, credit bureau information, payments data, public
records, firmographics, and statistical priors. This makes the share-of-wallet
problem partly a missing-data problem and partly a market-demand problem.

\subsection{Bank Deposit Demand and Competition}

Deposit wallet estimation cannot be reduced to generic cross-selling because
deposits have special economic features. Deposits are funding sources for banks,
liquidity services for customers, and relationship anchors. They also respond to
rates, fees, branch and digital convenience, safety perceptions, payment needs,
and switching frictions.

\citet{dick2008demand} estimates demand and welfare in the banking industry,
providing a structural demand perspective for bank services. \citet{adams2007who}
study competition among depository institutions, which is important for defining
the relevant competitive set. For an SME in a local market, the competitors may
not be all banks in the country; they may be nearby branch banks, digital banks,
specialized lenders, card issuers, fintech payment providers, and industry
specialists.

\citet{egan2017deposit} estimate deposit competition and financial fragility in
the US banking sector. Their work is relevant because deposit demand and bank
competition affect how much of a customer's balances can realistically be moved
by price or relationship offers. \citet{drechsler2017deposits} study the
deposits channel of monetary policy, emphasizing the role of deposit market
power and rate pass-through. For a wallet-size system, these papers motivate
models that include deposit rates, local competition, account type, market
concentration, and customer rate sensitivity.

A practical deposit model often needs a two-part structure:
\[
  \Pr(a_{if,\mathrm{dep}}>0\mid X_i)
\]
for deposit relationship incidence and
\[
  \mathbb{E}[\log a_{if,\mathrm{dep}}\mid a_{if,\mathrm{dep}}>0,X_i]
\]
for conditional deposit size. If competitor balances are unobserved, the model
may use market-level deposit data, firm revenue estimates, industry working
capital ratios, payments volume, and observed inflows/outflows as indirect
signals of total deposit wallet.

\subsection{Dynamic Deposit Demand}

A major weakness of the wallet-size analogy is the implicit constant-size
assumption. For deposits, the assumption is particularly fragile. An SME's
deposit balance is not only a product holding; it is a time-varying operating
balance. It moves with revenues, payroll, vendor payments, seasonality, business
growth, and the normal timing of cash inflows and outflows. A firm can need the
same banking product in two periods but require very different deposit balances.

The corporate liquidity literature gives a useful framework. \citet{opler1999cash}
study determinants of corporate cash holdings and emphasize transaction and
precautionary motives. \citet{almeida2004cash} develop the idea that financially
constrained firms save cash out of cash flows, while \citet{han2007precautionary}
study precautionary cash holdings under cash-flow risk. \citet{bates2009cash}
show that cash holdings of US firms increased over time and relate this to
riskier cash flows and changing firm characteristics. These papers are mostly
about corporate cash rather than bank deposit accounts, but for SME banking the
connection is useful: much of an SME's operating liquidity is expressed as bank
deposits or near-deposit balances.

Private-firm and smaller-firm evidence is especially relevant. \citet{mortal2020private}
compare cash holdings of private and public firms internationally.
\citet{poti2020privatecash} and \citet{lefebvre2021cash} study precautionary
cash motives in private and smaller firms. These papers support the view that a
deposit target should vary across industries, firm sizes, geographies, and
business cycles rather than being treated as a fixed customer attribute. The
more specialized question of deposit pricing and stress-period behavior is
discussed separately below.

This literature suggests that the focal object should be dynamic deposit demand,
not fixed deposit wallet size. A more appropriate target is
\[
  D_{it}^{*}
  =
  g(X_{it}, M_t, R_{it}, C_{it}, B_{it}, \eta_i) + \varepsilon_{it},
\]
where $D_{it}^{*}$ is desired liquidity or deposit balance, $X_{it}$ are
firm-level cash-flow and product variables, $M_t$ are market conditions, $R_{it}$
captures rates and fees, $C_{it}$ captures financing conditions, $B_{it}$
captures bank relationship and switching frictions, and $\eta_i$ is firm
heterogeneity. Observed deposits are then a partial allocation of this
desired liquidity across banks and instruments:
\[
  D_{ift} = s_{ift} D_{it}^{*},
\]
where $s_{ift}$ is the focal bank's dynamic share of deposit liquidity. In this
view, wallet estimation decomposes into two moving parts: the time-varying size
of the firm's liquidity demand and the focal bank's share of that demand.

\subsection{SME Relationship Lending and the Role of Transaction Accounts}

SME banking differs from mass retail banking because information asymmetry is
central. Relationship lending theory explains why small firms benefit from
long-lived bank relationships and why transaction data can be valuable for
credit decisions. \citet{petersen1994benefits} and
\citet{berger1995relationship} are foundational references on small-business
lending relationships and lines of credit.

Transaction accounts are particularly important. \citet{mester2007transactions}
study transaction accounts and loan monitoring. \citet{norden2010creditline}
examine credit-line usage, checking-account activity, and default risk. These
papers motivate a wallet model in which deposits and payments are not merely
products to be sold; they are information channels that can improve underwriting,
monitoring, and relationship profitability.

For an SME bank, this means a deposit offer may have value beyond spread income.
Moving an SME's operating account to the focal bank can reveal cash-flow
seasonality, payroll cycles, customer concentration, merchant activity, and early
distress signals. A product-propensity model should therefore be joined with a
relationship-value model.

\subsection{Causal Response and Uplift Modeling}

Cross-sell and wallet models estimate what an SME appears likely to buy or hold.
Marketing decisions require a causal estimand: what the SME will buy because of
the offer. A high-propensity customer may buy without being contacted; a
low-propensity customer may be expensive to convert; and an offer may shift
balances without increasing total relationship value. The relevant target is
incremental contribution, not raw propensity:
\[
  \tau_{ip}(x) =
  \mathbb{E}[Y_{ip}(1)-Y_{ip}(0)\mid X_i=x],
\]
where $Y_{ip}$ may be product adoption, deposit balance, risk-adjusted revenue,
or lifetime value.

The modern causal-ML literature provides tools for heterogeneous treatment
effects. \citet{radcliffe2007uplift} and \citet{lo2002true} discuss uplift and
true-lift modeling in marketing contexts. \citet{athey2016recursive} introduce
recursive partitioning for heterogeneous causal effects, and
\citet{wager2018causalforests} develop causal forests. These methods are useful
when a bank has randomized campaigns, quasi-random offer variation, or credible
experiments. Without such variation, a response model risks confusing selection
with treatment effect.

\subsection{Model Architectures Implied by the Literature}

The preceding literatures imply several architectures. These architectures are
not mutually exclusive. They answer different parts of the SME wallet problem,
and the main design question is how to combine them without confusing latent
need, observed product ownership, competitive share, and causal response.

\subsubsection*{Architecture 1: Latent Wallet Reconstruction}

The first architecture treats competitor wallet as a latent variable. A
hierarchical model can describe product incidence and size jointly:
\[
  z_i \sim p(z_i\mid X_i),
  \qquad
  y_{ibp} \sim \mathrm{Bernoulli}(\pi_{bp}(z_i,X_i)),
\]
\[
  \log a_{ibp}\mid y_{ibp}=1 \sim
  \mathcal{N}(\mu_{bp}(z_i,X_i),\sigma_{bp}^2),
\]
where $z_i$ represents latent firm scale, financial sophistication, liquidity
need, credit demand, and relationship preference. This architecture is closest
in spirit to latent trait and mixed-data factor models in the cross-selling
literature \citep{kamakura1991latent,kamakura2003database}. It also follows
the econometric logic of two-part or hurdle models: product incidence and
positive product size are distinct stochastic objects, so a bank should not
convert a product-propensity score into dollars without a conditional amount
model \citep{cragg1971limited}.

For implementation, a bank could start with a simpler generalized linear mixed
model or gradient-boosted two-part model and then move toward Bayesian
hierarchical or variational models as the data maturity improves. The output is
a posterior distribution for total wallet, competitor-held wallet, and product
amount conditional on ownership, not just a point estimate. If positive amounts
are modeled on the log scale, retransformation to dollar levels should be
validated rather than handled by naive exponentiation; smearing-style
corrections are a standard remedy \citep{duan1983smearing}.

\subsubsection*{Architecture 2: Product Graph and Next-Best-Action System}

The second architecture treats products as a graph or sequence. Product holdings
form a vector
\[
  \mathbf{y}_i = (y_{i1},\ldots,y_{iP}),
\]
and the model predicts the next product or bundle conditional on the current
state. This can be implemented using multivariate probit, discrete hazard
models, recurrent models, or gradient-boosted ranking models. The academic
foundation comes from cross-selling and NPTB models
\citep{knott2002nptb,li2005sequential}.

For SME banking, the graph should distinguish operating products from credit
products and owner/household products. For example, a business checking account
may lead to merchant services and payroll, while a line of credit may lead to
equipment finance or treasury management. A product graph is useful for sales
prioritization but insufficient for deposit sizing unless paired with balance
models. For the current project, this architecture is secondary; it is included
as background because it explains product correlations, not because assortment
or bundle choice is the main objective.

\subsubsection*{Architecture 3: Structural Deposit Demand and Competitive Response}

The third architecture models the SME's bank choice and deposit allocation as a
competitive demand problem. A simplified random-utility model is
\[
  U_{ibp} =
  \alpha_i r_{bp}
  - \phi_i F_{bp}
  + \gamma_i Q_{bp}
  - \kappa_i C_{ib}
  + \epsilon_{ibp},
\]
where $r_{bp}$ is the rate or yield, $F_{bp}$ fees, $Q_{bp}$ service quality,
and $C_{ib}$ switching or relationship cost. Structural banking papers motivate
including local competition, bank characteristics, and market power
\citep{dick2008demand,adams2007who,egan2017deposit,drechsler2017deposits}.

This architecture is most useful for pricing and deposit offers. It can answer
questions such as: how much incremental deposit balance should be expected if
the bank offers a rate premium, fee waiver, or improved onboarding support? Its
weakness is data intensity: credible estimation often needs variation in rates,
fees, markets, and competitor conditions.

\subsubsection*{Architecture 4: Dynamic Panel and Time-Series Deposit Demand}

Because deposits are balances observed over time, a bank should not rely only on
cross-sectional wallet models. A dynamic panel model can separate firm-specific
liquidity level, time-varying business conditions, market-wide stress, and
relationship share. A practical specification is
\[
  \log D_{ift}
  =
  \rho \log D_{if,t-1}
  + \beta' X_{it}
  + \theta' M_t
  + \psi' O_{ift}
  + \alpha_i
  + \delta_t
  + u_{ift},
\]
where $D_{ift}$ is the focal-bank deposit balance, $X_{it}$ includes observed
firm cash-flow and transaction features, $M_t$ includes macro and market
conditions, $O_{ift}$ includes the bank's offer and relationship terms,
$\alpha_i$ captures persistent firm heterogeneity, and $\delta_t$ captures
common shocks. This model can be extended with industry-by-time effects,
geography-by-time effects, state-space components, or regime-switching terms for
stress periods.

If the target is total deposit demand rather than focal-bank deposits, the bank
can model
\[
  \log D_{it}^{*}
  =
  \rho \log D_{i,t-1}^{*}
  + \beta' X_{it}
  + \theta' M_t
  + \alpha_i
  + \delta_t
  + \nu_{it}
\]
as a latent state and then model the focal bank's observed share:
\[
  \mathrm{logit}(s_{ift}) =
  \gamma' Z_{ift} + \mu_i + \xi_t.
\]
This separates the question ``how much liquidity does the SME want now?'' from
``what share of that liquidity sits with our bank?'' It also allows stress
periods to raise total precautionary balances while the focal bank's share may
rise, fall, or stay flat.

Time-series methods are useful, but they should usually be panel time-series
methods rather than a separate univariate model for each SME. Individual SME
series can be short, seasonal, lumpy, and affected by idiosyncratic events. A
hierarchical panel, state-space, or mixed-effects model can pool information
across similar firms while allowing customer-specific dynamics. For forecasting,
machine-learning sequence models can be useful, but they should be benchmarked
against transparent baselines such as seasonal naive forecasts, autoregressive
distributed-lag models, and fixed-effects panel regressions.

\subsubsection*{Architecture 5: Causal Uplift and Incremental Value}

The fifth architecture starts from experiments or quasi-experiments. The bank
randomizes or otherwise credibly varies offers and estimates heterogeneous
treatment effects. The targeting rule is then based on expected incremental
profit:
\[
  \mathrm{NPV}_{it}
  =
  \sum_p m_p \widehat{\Delta}_{ipt}
  - c_t
  - \lambda \widehat{\mathrm{Risk}}_{it},
\]
where $m_p$ is product margin or relationship value, $c_t$ is campaign and
incentive cost, and the risk term captures credit, liquidity, operational, and
conduct constraints. Uplift and causal-forest methods are natural candidates
\citep{lo2002true,radcliffe2007uplift,athey2016recursive,wager2018causalforests}.

This architecture is the most defensible for deciding whom to contact and what
to offer, because it targets incremental response rather than raw propensity.
Its prerequisite is disciplined experimentation and holdout design.

\subsubsection*{Architecture 6: Integrated Decision System}

The most complete architecture combines the preceding pieces:

\begin{enumerate}[leftmargin=*]
  \item Estimate total deposit wallet and competitor-held balances with a
  latent wallet model, and treat deposit wallet as time varying.
  \item Use product-adoption and product-graph models only as supporting
  signals for firm scale and relationship context.
  \item Estimate dynamic deposit demand using panel time-series models that
  include firm cash-flow, seasonality, macro conditions, rates, and local market
  conditions.
  \item Estimate deposit-specific response to rates, fees, service, and
  relationship actions with a demand model.
  \item Estimate the movable fraction of competitor-held balances, borrowing,
  spend, or assets from conversion events, account-transfer evidence,
  relationship migrations, and experiments.
  \item Estimate causal incremental lift from campaigns using randomized
  holdouts.
  \item Optimize offers under profitability, credit-risk, liquidity,
  relationship-manager capacity, fairness, and compliance constraints.
\end{enumerate}

This system should produce uncertainty intervals and decision diagnostics, not
only point scores. Relationship managers need to know whether a prospect has a
high expected deposit wallet, a high probability of product need, high
incremental responsiveness, or merely high observed similarity to existing
large customers.

\section{Data Sources and Feature Families}

The data problem is at least as important as the model class. No single public
or commercial data set observes the full SME wallet at competitor banks. The
bank therefore needs a layered data design: internal customer data identify the
focal-bank relationship, public data anchor local markets and stress
conditions, third-party and consented data provide partial external-wallet
signals, and experiments or conversion events identify move-over and
incrementality. The Federal Reserve data portal is useful for this program
because it provides official macro-financial, banking, interest-rate, lending,
and payments data \citep{fed2026data}. Its role, however, is mostly to supply
market conditions, rate regimes, credit conditions, payment trends, and stress
context. It does not directly identify an individual SME's deposits at
competitor banks.

\subsection{Required Data Roles}

The literature implies that an SME wallet-size model should combine several
types of data, each with a different identification role:

\begin{enumerate}[leftmargin=*]
  \item \textbf{Internal relationship data.} Observed deposits, loan balances,
  card spend, merchant transactions, treasury services, tenure, service events,
  relationship-manager notes, customer-service history, onboarding channel,
  historical offers, quoted rates, booked rates, fee waivers, and exceptions.
  These data define the focal-bank labels and the information set available at
  decision time.
  \item \textbf{Transaction and operating-flow signals.} ACH, wire, payroll,
  merchant, check, card, lockbox, remote-deposit, and cash-flow patterns can
  proxy firm scale, seasonality, payroll intensity, supplier/customer networks,
  and external-bank relationships. Inflow and outflow volatility, payroll
  coverage, tax-payment timing, merchant settlement volume, and counterparty
  concentration are especially important for estimating total liquidity demand.
  \item \textbf{Firmographics and local business structure.} Industry,
  geography, firm age, revenue band, payroll size, employee count, legal form,
  ownership, seasonality, and local establishment density. Public sources such
  as County Business Patterns, Business Dynamics Statistics, Business Formation
  Statistics, and QCEW help calibrate industry-by-geography opportunity and
  business-cycle exposure \citep{census2026cbp,census2026bds,census2026bfs,bls2026qcew}.
  \item \textbf{Credit and public records.} Liens, UCC filings, commercial real
  estate loans, commercial credit bureau data, SBA lending records, property
  records, bankruptcy events, and public registrations where legally and
  contractually usable. These data are useful for product incidence and size,
  especially loans, lines, collateralized credit, commercial real estate
  exposure, and distress events.
  \item \textbf{Market, competitor, and branch data.} Local branch density,
  branch-level deposit market share, competitor footprint, advertised rates,
  fee schedules, product availability, digital-bank and fintech presence, and
  local competitive intensity. FDIC Summary of Deposits is a natural public
  anchor for branch-level deposits \citep{fdic2026sod}. The limitation is that
  such aggregates usually mix consumer, small-business, commercial, and
  wholesale deposits, so they should be used as noisy market-level anchors
  rather than direct SME wallet labels.
  \item \textbf{Rate and stress context.} Federal Reserve H.8 gives aggregate
  commercial-bank balance-sheet conditions, H.15 gives selected interest rates,
  and the Senior Loan Officer Opinion Survey gives credit-standard and
  loan-demand context
  \citep{fed2026h8,fed2026h15,fed2026sloos}. Call Reports from FDIC or FFIEC
  provide bank-level balance-sheet and funding information
  \citep{fdic2026callreports,ffiec2026cdr}. These data support market regimes,
  competitor funding pressure, stress scenarios, and deposit-pricing variables.
  \item \textbf{Payments and commerce benchmarks.} Federal Reserve payments
  data, card-network benchmarks, merchant-acquirer data, payroll-provider
  signals, and industry payment benchmarks help translate operating activity
  into expected deposit balances and card or merchant-wallet opportunity
  \citep{fed2026frps}. They are usually aggregate or vendor-provided, so they
  need reconciliation against internal transaction data.
  \item \textbf{Validation anchors.} Consented account aggregation,
  customer-provided statements, treasury onboarding documents, balance-transfer
  records, payoff letters, direct-deposit or ACH rerouting, account-closure
  evidence, M\&A onboarding files, and post-conversion consolidated balances are
  the most valuable data for validating competitor wallet estimates. Without
  these anchors, the model may rank prospects but should not claim precise
  competitor-dollar estimates.
  \item \textbf{Campaign, quote, and experiment data.} Offer eligibility,
  contact channel, price, incentive, quoted and unquoted terms,
  relationship-manager activity, randomized holdouts, quasi-random rate cells,
  and rejected offers. These data identify price sensitivity, causal uplift,
  move-over, cannibalization, and runoff.
\end{enumerate}

\subsection{How Specific Public Data Help}

Table~\ref{tab:data_sources} summarizes the most useful public and semi-public
sources. The important point is not that these data replace bank data. They
provide anchors, controls, market states, and validation benchmarks that make
the internal model less circular.

\begin{table}[htbp]
\centering
\footnotesize
\begin{tabular}{@{}
  >{\raggedright\arraybackslash}p{0.23\linewidth}
  >{\raggedright\arraybackslash}p{0.24\linewidth}
  >{\raggedright\arraybackslash}p{0.25\linewidth}
  >{\raggedright\arraybackslash}p{0.18\linewidth}
@{}}
\toprule
Source & Useful variables & Role in the model & Main limitation \\
\midrule
Federal Reserve data portal & H.8, H.15, SLOOS, payments data, business-finance
releases & Macro state, rate regime, credit conditions, payments context, stress
scenario variables & Mostly aggregate, not SME wallet labels \\
FDIC Summary of Deposits & Branch deposits, institution and geography footprint
& Local deposit market share and branch competition anchor & Mixes consumer,
SME, commercial, and wholesale deposits \\
FDIC/FFIEC Call Reports & Bank-level deposits, loans, funding mix, uninsured
deposit measures where reported & Competitor funding pressure, safety, pricing,
and stress controls & Bank-level, not customer-level \\
Census CBP/BDS/BFS and BLS QCEW & Establishments, firm births, employment,
wages, industry geography & Local SME universe, market size, growth, seasonality
and industry stress & Establishment-level aggregates, not bank relationships \\
CRA, HMDA, and Section 1071 data & Small-business and mortgage lending
geography, where available & Credit demand and market access context; future
validation of small-business credit coverage
\citep{ffiec2026cra,cfpb2026hmda,cfpb2023section1071} & Product scope and
reporting definitions differ from deposit wallet \\
Commercial/vendor data & Firmographics, credit bureau, card/merchant/payroll,
web and business registry data & Product incidence, size priors, external-wallet
signals, prospect universe & Cost, coverage bias, consent, privacy, and model
risk review \\
Consented or conversion data & External account balances, statement imports,
ACH migration, payoff and transfer records & Direct validation of competitor
wallet, move-over, cannibalization, and runoff & Sparse and selected sample \\
\bottomrule
\end{tabular}
\caption{Data sources for SME wallet estimation and their identification roles.}
\label{tab:data_sources}
\end{table}

\subsection{Data Architecture for the Recommended Model}

For the recommended architecture, data should be organized around estimands
rather than around source systems.

\paragraph{Total deposit demand.}
Estimate $D_{it}^{*}$ from internal transaction flows, firmographics, seasonal
patterns, local industry conditions, rate regimes, credit conditions, and
external business-density benchmarks. Federal Reserve rate and credit data enter
as $M_t$ and stress-state variables; Census and BLS data enter as
industry-geography controls; transaction flows identify firm-specific operating
liquidity.

\paragraph{Competitor share and local substitution.}
Estimate $s_{ifpt}$ using focal-bank balances, branch footprint, competitor
footprint, advertised rates, service attributes, branch-level market share, and
bank-level balance-sheet controls. FDIC Summary of Deposits is useful here as a
noisy local-market share anchor, while Call Reports and H.8/H.15 help measure
funding and rate environments. These data cannot by themselves separate SME
from consumer or wholesale share; the model must treat them as aggregate
measurement equations.

\paragraph{Product incidence and size.}
Estimate product incidence and conditional dollar size from internal products,
credit bureau and public-record features, SBA or public lending records where
available, firmographics, and transaction proxies. For loans and credit lines,
liens, UCC filings, property records, and credit bureau data are more direct
than Fed aggregate data. For cards and merchant products, merchant settlement
and payment-flow benchmarks are more relevant.

\paragraph{Move-over, cannibalization, and runoff.}
Estimate the movable fraction using data that reveal origin and displacement:
account aggregation, conversion files, balance-transfer and payoff records,
ACH/direct-deposit migration, customer onboarding statements, relationship-
manager disposition codes, and randomized consolidation offers. This is the
hardest part of the data problem. Public data can benchmark market size, but
they cannot identify whether a given balance came from a competitor, from new
liquidity demand, or from cannibalization of existing focal-bank balances.

\paragraph{Pricing and stress.}
Estimate pricing response from quoted rates, booked rates, fee waivers,
rate-sheet changes, market rates, competitor advertised rates, and randomized
or quasi-random offer variation. H.15, H.8, SLOOS, Call Reports, and FDIC market
data provide the market and stress states. Internal balance decay and
post-stress runoff are still required to value the deposits.

\subsection{Practical Data Priorities}

The most valuable next data investments are:
\begin{enumerate}[leftmargin=*]
  \item Build an account-level deposit and product panel with decision-time
  snapshots, so the model never uses post-offer information in scoring.
  \item Create a transaction-feature mart for cash-flow volatility, payroll
  intensity, merchant settlement, ACH/wire counterparties, tax timing, and
  inferred external-bank activity.
  \item Geocode customers, branches, and competitors; join FDIC Summary of
  Deposits, Census/BLS local business data, and competitor footprints.
  \item Maintain a quote and offer ledger with offered rate, booked rate,
  rejected terms, fee waivers, exceptions, relationship-manager action, and
  eligibility rules.
  \item Create validation samples through consented account aggregation,
  onboarding statement capture, conversion events, and relationship-manager
  migration codes.
  \item Run small randomized or quasi-random pricing and consolidation tests to
  identify elasticity, move-over, cannibalization, and runoff.
\end{enumerate}

Data governance is a first-order issue. A US bank must ensure that features,
target definitions, offer decisions, and model monitoring comply with fair
lending, privacy, consumer/commercial data-use restrictions, model-risk
management, and customer consent requirements. The academic literature reviewed
here does not by itself settle those compliance requirements.

For deployment, the modeling program should be governed as a bank model-risk
and customer-decisioning system, not only as a marketing model. SR 11-7 requires
conceptual soundness, ongoing monitoring, and independent validation for models
used in bank decisioning \citep{fed2011sr1107}. Third-party data and fintech
components should be reviewed under third-party risk-management expectations
\citep{fed2023sr2304}. If the model is used in connection with small-business
credit decisions or covered data collection, implementation should also be
reviewed against the CFPB small-business lending rulemaking under Section 1071
\citep{cfpb2023section1071}.

\section{Recommended Approach}

The recommended approach is dynamic demand plus discrete choice. It begins by
abandoning the idea that each SME has a fixed deposit wallet size. The bank
should instead model two linked objects:
\[
  \text{Observed focal-bank deposits}_{ift}
  =
  \text{Total SME liquidity demand}_{it}
  \times
  \text{Focal-bank deposit share}_{ift}.
\]
The first object is a dynamic balance-sheet and cash-flow problem. The second
object is a bank-choice and deposit-share problem. A third object is the
deposit amount that can move from competitors, which is generally
smaller than the competitor-held amount. This decomposition is the central
modeling recommendation of this survey.

\subsection{Layer 1: Dynamic SME Liquidity Demand}

The first layer estimates the SME's desired liquidity or total deposit demand:
\[
  D_{it}^{*}
  =
  g(X_{it}, M_t, R_{it}, C_{it}, B_{it}, \eta_i) + \varepsilon_{it}.
\]
Here $D_{it}^{*}$ is the firm's desired liquid balance, $X_{it}$ includes
firm-level cash-flow and transaction variables, $M_t$ includes macro and market
conditions, $R_{it}$ includes relevant rates and fees, $C_{it}$ includes
financing conditions, $B_{it}$ includes relationship variables, and $\eta_i$
captures persistent firm heterogeneity. This layer should be estimated as a
panel time-series, state-space, or hierarchical dynamic model. It should include
seasonality, industry-by-time effects, geography-by-time effects, local market
conditions, and relationship tenure. Pricing and stress extensions can be added
after the baseline model is stable.

This layer answers questions such as: how much liquidity does the SME need now,
how does that need vary with business scale and seasonality, and how much of a
balance increase reflects normal cash-flow timing rather than wallet movement
from a competitor?

\subsection{Layer 2: Discrete Choice for Bank and Deposit Share}

The second layer estimates how the SME allocates that liquidity and product
demand across banks and outside options. This is where discrete
choice demand estimation is useful. The alternative set must be declared before
estimation. In the deposit use case, let $\mathcal{J}_{imt}^{B}$ denote modeled
bank-relationship alternatives, let $\mathcal{J}_{imt}^{BP}$ denote
bank-product alternatives when products are modeled separately, let
$j=\mathrm{res}$ denote residual unmodeled-bank balances, and let
$j=0$ denote true nonbank or outside liquidity such as money-market funds,
cash-management vehicles, or no active movement. These categories should not
be pooled unless the reporting question only requires total non-focal
balances.

A simplified utility specification for SME $i$ choosing alternative $j$ in
market $m$ and time $t$ is
\[
  U_{ijmt}
  =
  \beta_i' A_{jmt}
  - \alpha_i P_{ijmt}
  + \xi_{jmt}
  + \epsilon_{ijmt},
\]
where $A_{jmt}$ are observed bank/product attributes, $P_{ijmt}$ is the
SME-specific price or offer term, $\xi_{jmt}$ is an unobserved product-market
demand shock, and $\epsilon_{ijmt}$ is idiosyncratic taste. Alternatives may be
defined as competing operating-account relationships or bank-product bundles,
but the same model should not switch among these definitions without changing
the share normalization and diversion interpretation.

\citet{ma2024discrete} is useful here because financial products often involve
individual-specific or relationship-specific prices. In SME banking, effective
prices may include negotiated deposit rates, treasury fees,
minimum-balance requirements, fee waivers, and relationship pricing. If these
prices vary systematically with firm size, risk, income, relationship tenure,
or financing conditions, a standard BLP model that treats all
customers in a market as facing the same product price can misestimate price
sensitivity. Ma's two-step framework---first estimating individual-specific
prices, then using a modified BLP-style demand model with selectivity
correction---is therefore relevant for estimating SME choice among financial
product alternatives.

Because \citet{ma2024discrete} is an unpublished working paper, it should be
treated as a useful template for handling individual-specific financial-product
prices rather than as settled banking practice. A bank should validate the
two-step price and demand estimates against its own quote, booking, and
campaign data before using them for pricing or targeting decisions.

For plain retail deposit products where all customers in a market truly face
the same posted rate, a traditional BLP-style model may be sufficient. For SME
relationship banking, however, prices and service terms are often negotiated or
relationship-specific, so the individual-specific-price issue is likely
material.

\subsection{Layer 3: Product Size, Addressable Wallet, and Move-Over Amount}

The third layer estimates dollars, balances, spend, commitments, or assets
conditional on product need. For every product $p$, the model should separate
incidence from severity:
\[
\begin{aligned}
  r_{ipt}^{*} &= 1\{W_{ipt}^{*}>0\},\\
  \Pr(r_{ipt}^{*}=1\mid h_{it},X_{it},M_t) &= \pi_p(h_{it},X_{it},M_t),\\
  \log W_{ipt}^{*}\mid r_{ipt}^{*}=1
    &\sim F_p(\mu_p(h_{it},X_{it},M_t),\sigma_p^2).
\end{aligned}
\]
This two-part structure follows the limited-dependent-variable logic of
\citet{cragg1971limited}. It prevents a common mistake in wallet sizing:
treating a product-propensity score as if it were a dollar balance. For
products with highly skewed positive amounts, the model should output means,
quantiles, and interval forecasts. If the severity model is fit in logs, dollar
forecasts should use validated retransformation, such as smearing corrections,
rather than naive exponentiation \citep{duan1983smearing}.

Given a draw of total product wallet and focal-bank share, the model produces
the competitor-held amount
\[
  A_{ipt}^{-f(s)}
  =
  \max\{W_{ipt}^{*(s)}-a_{ifpt}^{\mathrm{obs}},0\}.
\]
If the bank needs bank-specific competitor estimates, this amount can be
allocated across competitors using the discrete-choice shares:
\[
  a_{ibpt}^{(s)}
  =
  s_{ibpt}^{(s)} A_{ipt}^{-f(s)}
  \bigg/
  \sum_{c\neq f}s_{icpt}^{(s)},
  \qquad b\neq f.
\]
This allocation is only over modeled banks. If important banks are not modeled,
include a residual unmodeled-bank category in the denominator and report it
separately. A true nonbank outside option should be allocated dollars only when
the wallet definition includes nonbank liquidity; otherwise it belongs in the
share model but not in bank-specific competitor balances.

The move-over amount is a separate estimand. For offer $o$ and horizon $H$,
write
\[
  \Delta a_{iop}^{\mathrm{move}}(H)
  =
  M_{iop}(H)A_{ipt}^{-f},
  \qquad 0\leq M_{iop}(H)\leq 1.
\]
The fractional response $M_{iop}(H)$ can be modeled with a fractional-logit,
beta-regression, hierarchical Bayesian, or causal-forest layer using offer
terms, switching costs, relationship depth, competitor attributes, price gaps,
and product frictions. It must be estimated from data that observe or identify
movement: conversion events, balance-transfer or payoff records, account
aggregation before and after onboarding, ACH or wire migration patterns,
relationship migrations, merger/acquisition integrations, or randomized offer
cells.

Observed focal-bank balance growth is not itself the move-over amount. The
accounting identity should be
\[
  \Delta a_{ifp,t:t+H}^{\mathrm{obs}}(o)
  =
  \Delta a_{iop}^{\mathrm{move}}(H)
  +\Delta a_{iop}^{\mathrm{expand}}(H)
  -\Delta a_{iop}^{\mathrm{cann}}(H)
  +u_{iopH},
\]
where $\Delta a^{\mathrm{move}}$ is migrated competitor wallet,
$\Delta a^{\mathrm{expand}}$ is expansion of the SME's total demand for product
$p$, and $\Delta a^{\mathrm{cann}}$ is cannibalization of existing focal-bank
balances or organic growth. Thus a deposit initiative can create large observed
growth because the firm expanded, because it moved competitor balances, or
because it shifted value from existing focal-bank balances. These components
have different profitability, funding, liquidity, and credit implications.

\subsection{Layer 4: Counterfactual Offer and Wallet-Share Simulation}

Once the demand, choice, and amount layers are estimated, the bank can run
counterfactuals. For a candidate deposit action $o$,
\[
  \widehat{D}_{ifo,t+1}
  =
  \widehat{D}^{*}_{i,t+1}
  \times
  \widehat{s}_{ifo,t+1},
\]
where $\widehat{D}^{*}_{i,t+1}$ comes from the dynamic liquidity-demand layer
and $\widehat{s}_{ifo,t+1}$ comes from the discrete-choice share layer. This
lets the bank distinguish three cases that look similar in a static wallet
score:

\begin{enumerate}[leftmargin=*]
  \item The SME has high total liquidity demand, but the focal bank has low
  share because the bank's price, service, or relationship position is
  uncompetitive.
  \item The SME has low current liquidity demand, but high future demand under
  seasonal growth.
  \item The SME has high observed balances at the focal bank, but low
  incremental response to marketing because it would have stayed anyway.
\end{enumerate}

The discrete-choice layer can also quantify willingness to pay for attributes:
branch access, relationship-manager quality, digital onboarding, faster credit
decisions, integrated treasury tools, deposit-rate premium, fee waivers, or
onboarding support. These willingness-to-pay estimates are useful for
evaluating deposit capture actions rather than only ranking prospects.

\subsection{Layer 5: Causal Uplift and Risk-Adjusted Decisioning}

The dynamic demand and discrete-choice layers estimate demand and substitution.
They do not by themselves prove that a marketing action causes incremental
balances. The final layer should use randomized holdouts or credible
quasi-experiments to estimate incremental response:
\[
  \tau_{io}^{Y}
  =
  \mathbb{E}[Y_i(o)-Y_i(0)\mid \mathcal{I}_{it}].
\]
The superscript matters: $\tau_{io}^{D}$ is a deposit-balance contrast,
$\tau_{io}^{\mathrm{move}}$ is migrated competitor wallet, and
$\tau_{io}^{V}$ is risk-adjusted value. They need not have the same sign or
the same identification support.
The recommended targeting criterion is risk-adjusted incremental value:
\[
  \mathrm{NPV}_{io}
  =
  m_D \widehat{\Delta D}_{io}
  + \sum_p m_p \widehat{\Delta a}_{iop}
  - c_o
  - \lambda \widehat{\mathrm{Risk}}_{io},
\]
where $\widehat{\Delta D}_{io}$ is incremental deposit balance, $m_D$ is the
deposit margin or funding value, $m_p$ are product margins, $c_o$ is offer cost,
and the risk term captures credit, liquidity, conduct, and compliance
constraints.

Thus the recommended system is:
\[
\resizebox{0.96\linewidth}{!}{$\displaystyle
  \boxed{
  \text{Dynamic liquidity demand}
  \;+\;
  \text{discrete choice deposit share}
  \;+\;
  \text{product amount and move-over}
  \;+\;
  \text{causal uplift}
  \;+\;
  \text{risk-adjusted optimization}.}
$}
\]
This framework keeps the core economic distinction clear: deposit balances are
dynamic demand, while wallet capture is a competitive choice and response
problem.

\subsection{Summary of the Recommended Model}

The recommended model combines the surveyed architectures into one
generative-and-decision system. This subsection gives the compact version; the
next two subsections state the estimands, estimator interfaces, and validation
gates. The key notation is:
\begin{center}
\footnotesize
\begin{tabular}{@{}p{0.18\linewidth}p{0.74\linewidth}@{}}
\toprule
Symbol & Meaning \\
\midrule
$W_{ipt}^{*}$ & Latent total wallet for product $p$ \\
$A_{ipt}^{-f}$ & Competitor-bank-held amount, excluding true nonbank outside liquidity unless stated \\
$D_{it}^{*}$ & Latent total SME deposit demand \\
$s_{ifpt}(o)$ & Focal-bank share of product $p$ under offer $o$ \\
$M_{iop}(H)$ & Structural fraction of competitor-held product $p$ that moves under offer $o$ \\
$\mu_{iop}(H)$ & Expected moved dollars \\
$\tau_{io}^{Y}$ & Causal contrast for explicitly named outcome $Y$ \\
$J_{io}(H)$ & Risk-adjusted incremental relationship value \\
\bottomrule
\end{tabular}
\end{center}
Let $h_{it}$ denote a latent SME state containing
firm scale, liquidity need, credit demand, relationship preference, price
sensitivity, and switching cost. The state evolves over time:
\[
  h_{it}
  =
  F(h_{i,t-1},X_{it},M_t,L_{it},\omega_{it}),
\]
where $X_{it}$ are firm signals, $M_t$ are market conditions, $L_{it}$ captures
financing conditions and liquidity constraints, and $\omega_{it}$ is an
innovation. Deposit need and total product wallet are generated from this state:
\[
  W_{ipt}^{*} \sim p(W_{ipt}^{*}\mid h_{it},X_{it},M_t),
  \qquad
  y_{ipt}^{*} \sim p(y_{ipt}^{*}\mid h_{it},W_{ipt}^{*}).
\]
For deposits, the product-wallet equation should be specialized as dynamic
liquidity demand:
\[
  D_{it}^{*}(M_t)
  =
  g(h_{it},\text{cash-flow risk}_{it},M_t,\text{relationship}_{it}).
\]
The allocation of that demand across the focal bank, competitors, and outside
options is then governed by a discrete-choice share model. Let
$\mathcal{B}_{mt}^{\mathrm{mod}}$ be the modeled bank set, let
$j=\mathrm{res}$ denote residual unmodeled-bank balances, and let $j=0$ denote
true nonbank outside liquidity. Then
\[
  s_{ibpt}(o)
  =
  q_{bp}(A_{bpt},P_{ibpt}(o,M_t),B_{ibt},h_{it},M_t,\xi_{bpt}),
  \qquad
  a_{ibpt}(o,M_t)=s_{ibpt}(o,M_t)W_{ipt}^{*}(M_t).
\]
Thus deposit pricing and stress enter through both total demand
$D_{it}^{*}(M_t)$ and captured share $s_{ifpt}(o,M_t)$.
The competitor-held and addressable components are
\[
  A_{ipt}^{-f}(0)=
  \sum_{b\in\mathcal{B}_{mt}^{\mathrm{mod}},\,b\neq f}a_{ibpt}(0)
  +a_{i,\mathrm{res},pt}(0),
  \qquad
  A_{ipo}^{\mathrm{addr}}(H)=M_{ipo}(H)A_{ipt}^{-f}(0).
\]
The nonbank outside option $j=0$ is included in $A_{ipt}^{-f}$ only if the
business definition of wallet includes nonbank liquidity. Otherwise it affects
choice normalization but is reported separately from competitor-bank wallet.
The causal move-over amount satisfies the accounting decomposition
\[
  \Delta a_{ifp,t:t+H}(o)
  =
  \Delta a_{iop}^{\mathrm{move}}(H)
  +\Delta a_{iop}^{\mathrm{expand}}(H)
  -\Delta a_{iop}^{\mathrm{cann}}(H),
\]
where only the first term is migrated competitor wallet. The second term is
growth in total product demand, and the third term is cannibalization of
existing focal-bank balances.
Finally, campaign experiments or quasi-experiments estimate the incremental
effect of taking action $o$ on explicitly named outcomes. For risk-adjusted
relationship value,
\[
  \tau_{io}^{V}
  =
  \mathbb{E}\{V_i(o)-V_i(0)\mid h_{it},X_{it}\},
\]
where $V_i(o)$ is risk-adjusted relationship value. This is distinct from
$\tau_{io}^{D}$ for deposit balances and $\tau_{io}^{\mathrm{move}}$ for
migrated competitor dollars. The decision rule is to
choose offers that maximize expected incremental value subject to eligibility,
fairness, liquidity, risk, capital, and relationship-manager capacity
constraints.

This combined model is more complete than any single surveyed architecture.
Latent wallet reconstruction fills the missing competitor-balance problem.
Dynamic liquidity demand prevents the deposit target from being treated as a
fixed stock. Two-part amount modeling estimates the size of each product
conditional on ownership. Discrete choice supplies the substitution and pricing
mechanism needed for counterfactual deposit actions. Section 6 extends this
same demand-and-share decomposition to stress-condition pricing, rate
sensitivity, and runoff. Move-over modeling distinguishes
competitor balances that exist from balances that can be captured. Causal
uplift prevents the bank from confusing high propensity with incremental
response. Relationship-lending research supplies the reason to optimize the
deposit relationship rather than the deposit margin in isolation.

\subsection{Why This Approach Is Preferable: Estimands and Losses}

The practical decision is not to estimate a scalar ``wallet size.'' It is to
choose an offer $o$ from a feasible set $\mathcal{O}_{it}$ using information
$\mathcal{I}_{it}$ available at decision time. A rigorous target is the
conditional incremental relationship value over horizon $H$:
\[
  J_{io}(H)
  =
  \mathbb{E}\!\left[
    V_{i,t+H}(o)-V_{i,t+H}(0)
    \mid \mathcal{I}_{it},\, o\in\mathcal{O}_{it}
  \right],
\]
where $V$ includes deposit funding value, product margin, servicing cost,
credit loss, liquidity stability, incentive cost, and compliance constraints.
The policy target is
\[
  \pi^{*}(\mathcal{I}_{it})
  =
  \arg\max_{o\in\mathcal{O}_{it}}
  J_{io}(H)
  \quad
  \text{subject to risk, fairness, capacity, and eligibility constraints.}
\]
A score is decision-relevant only to the extent that it is monotone in
$J_{io}(H)$ for feasible offers. A fixed wallet-size score generally is not,
because it omits dynamic liquidity demand, competitive substitution, treatment
response, product-specific amount distributions, movable competitor wallet, and
risk-adjusted value.

For deposits, the decomposition makes the failure precise. Let
$D_{ift}(o)=s_{ift}(o)D_{it}^{*}(o)$ denote the focal-bank deposit balance under
offer $o$, where $D_{it}^{*}$ is total SME liquidity demand and $s_{ift}$ is the
focal bank's share. Relative to no action,
\[
\begin{aligned}
  \Delta D_{if,t+H}(o)
  &=
  D_{i,t+H}^{*}(0)\Delta s_{if,t+H}(o)
  + s_{if,t+H}(0)\Delta D_{i,t+H}^{*}(o) \\
  &\quad
  + \Delta s_{if,t+H}(o)\Delta D_{i,t+H}^{*}(o),
\end{aligned}
\]
where $\Delta s_{if,t+H}(o)=s_{if,t+H}(o)-s_{if,t+H}(0)$ and
$\Delta D_{i,t+H}^{*}(o)=D_{i,t+H}^{*}(o)-D_{i,t+H}^{*}(0)$. A static wallet
model may forecast net focal-bank balance changes, but it does not separately
identify these channels. In particular, it cannot tell whether the observed
change came from share capture, changes in the firm's total liquidity demand,
the interaction of the two, or causal incrementality of the offer without
additional structure or experimental variation.

For non-deposit products the same issue appears as an incidence-size-movement
decomposition. Let $r_{ipt}^{*}=1\{W_{ipt}^{*}>0\}$ and let $M_{iop}(H)$ denote
the movable fraction of competitor-held product $p$. The gross move-over target
is
\[
  G_{iop}(H)
  =
  \mathbb{E}\!\left[
  1\{r_{ipt}^{*}=1\}A_{ipt}^{-f}M_{iop}(H)
  \mid \mathcal{I}_{it},o
  \right].
\]
This formula shows why three models are required: product incidence, conditional
amount, and offer-specific movable fraction. Estimating only the first factor
answers ``who may need the product,'' not ``how many dollars are at competitors''
or ``how many dollars are likely to move.'' A factorized approximation can be
used for transparency, but it should be validated because competitor-held
amount and movable fraction are likely correlated through firm size, switching
costs, relationship depth, and competitor lock-in.

The fixed-wallet approach is valid only under strong assumptions:
\begin{enumerate}[leftmargin=*]
  \item $D_{i,t+H}^{*}(o)=D_{it}^{*}$, so the firm's liquidity demand is fixed
  over the horizon and unaffected by market conditions or the offer.
  \item $s_{if,t+H}(o)$ depends only on an observed propensity score, not on
  unobserved competitor attributes, negotiated prices, switching frictions, or
  relationship-manager actions.
  \item Product incidence is enough to determine dollar size, so conditional
  amount distributions and retransformation uncertainty can be ignored.
  \item The movable fraction $M_{ipo}(H)$ is either one or a known constant,
  so all competitor-held balances are treated as equally addressable.
  \item The offer assignment mechanism is ignorable:
  $Y_i(o)\perp O_i\mid \mathcal{I}_{it}$, or there is randomized or
  quasi-random variation that identifies the counterfactual.
  \item Product margins, balance stability, credit risk, and servicing cost are
  either constant across SMEs or irrelevant to the decision.
  \item Moving deposits does not substitute away from other profitable products
  or cannibalize existing focal-bank balances.
\end{enumerate}
The literatures surveyed above weaken these assumptions: corporate liquidity
research rejects time-invariant cash demand
\citep{opler1999cash,almeida2004cash,bates2009cash,acharya2020dash};
structural demand models treat price and product attributes as drivers of share
\citep{berry1995automobile,dick2008demand}; causal-response methods separate
propensity from treatment effect
\citep{lo2002true,athey2016recursive,wager2018causalforests}; and
relationship-lending research implies that deposits, monitoring, and credit
products jointly affect value
\citep{petersen1994benefits,berger1995relationship,mester2007transactions,norden2010creditline}.

\begin{table}[htbp]
\centering
\footnotesize
\begin{tabular}{@{}
  >{\raggedright\arraybackslash}p{0.18\linewidth}
  >{\raggedright\arraybackslash}p{0.22\linewidth}
  >{\raggedright\arraybackslash}p{0.26\linewidth}
  >{\raggedright\arraybackslash}p{0.22\linewidth}
@{}}
\toprule
Model class & Estimand it targets & Failure mode for SME deposits & Role in recommended system \\
\midrule
Fixed wallet score & Rank or point estimate of total balances & Assumes liquidity demand and wallet are stable; confounds size and share & Baseline comparator only \\
Cross-sell/NPTB & Product incidence or next-product probability & Predicts need, not balance size or causal response & Product graph and incidence prior \\
Two-part amount model & $W_{ipt}^{*}$, $A_{ipt}^{-f}$, and conditional product size & Needs anchors and careful dollar retransformation & Product-size and competitor-wallet layer \\
Dynamic deposit model & $D_{it}^{*}$ or $D_{ift}$ over time & Forecasts balances but not competitive substitution by itself & Liquidity-demand layer \\
Discrete choice & $s_{ibpt}(o)$ and substitution elasticities & Needs prices, attributes, and identification of price effects & Deposit-share and offer-simulation layer \\
Move-over/uplift & $M_{ipo}(H)$, $\mu_{iop}(H)$, $J_{io}(H)$, or $\tau_{io}^{Y}$ under an action & Needs migration evidence and randomized or credible quasi-random variation & Campaign-response and policy layer \\
\bottomrule
\end{tabular}
\end{table}

\subsection{How the Literature Components Fit Together: Estimation and Interfaces}

The architecture should be written as a set of estimators with explicit
interfaces. Let
\[
  \mathcal{I}_{it}
  =
  \sigma\{X_{i,1:t},M_{1:t},A_{1:t},P_{i,1:t}^{\mathrm{obs}},
  a_{if,1:t}^{\mathrm{obs}},Y_{i,1:t}^{\mathrm{obs}},O_{i,1:t-1}\}
\]
denote the information available before the period-$t$ decision. The system
must not use information realized after the offer when estimating a decision
score. Let $\mathcal{A}_{it}$ denote anchoring data, such as account-aggregation
records, post-conversion balances, customer self-reports, bureau or public
records, branch-level market shares, and market-cell totals. These anchors are
essential because internal balances alone do not identify competitor balances.
Branch-level market share is informative but mixed: it usually combines
consumer, SME, commercial, and wholesale deposits. It should therefore enter as
an aggregate constraint or noisy measurement, not as a direct SME label.

\paragraph{Module A: latent wallet and state filtering.}
The target is the filtering distribution
\[
  F_{it}^{W,h}
  =
  p(h_{it},\{W_{ipt}^{*},A_{ipt}^{-f}\}_p,D_{it}^{*}
  \mid \mathcal{I}_{it},\mathcal{A}_{it}),
\]
not a point estimate. A useful measurement system is
\[
\begin{aligned}
  r_{ipt}^{*} &= 1\{W_{ipt}^{*}>0\},\\
  \log W_{ipt}^{*}\mid r_{ipt}^{*}=1
    &= \mu_p(h_{it},X_{it},M_t)+u_{ipt},\\
  a_{ifpt}^{\mathrm{obs}}
    &= s_{ifpt}W_{ipt}^{*}+\varepsilon_{ifpt}^{a},\\
  A_{ipt}^{\mathrm{anchor}}
    &= W_{ipt}^{*}+\varepsilon_{ipt}^{A}
    \quad \text{when an anchor is observed,}\\
  S_{fmt}^{\mathrm{branch}}
    &= \sum_g \omega_{gmt}S_{fgmt}+\varepsilon_{fmt}^{S}.
\end{aligned}
\]
The anchor equations are what provide level information for the latent wallet.
Without individual anchors, branch-market aggregates, or credible market-cell
restrictions, the model may rank SMEs but cannot claim point identification of
total competitor wallet. Mixed branch aggregates should be treated as noisy
constraints, not as direct SME wallet observations. The
estimator can be maximum likelihood with EM, Bayesian data augmentation, or a
mixed-data factor model
\citep{kamakura1991latent,kamakura2003database,dempster1977em}. For product
amounts, the incidence and positive-size equations should be estimated and
validated separately, as in two-part limited dependent variable models
\citep{cragg1971limited}. Its required output is a set of draws
\[
  \{h_{it}^{(s)},W_{ipt}^{*(s)},A_{ipt}^{-f(s)},D_{it}^{*(s)}\}_{s=1}^{S},
\]
which preserves wallet uncertainty for downstream choice and policy
simulation.

\paragraph{Module B: dynamic deposit demand.}
For deposits, write $d_{it}^{*}=\log D_{it}^{*}$. A minimal panel state equation
is
\[
  d_{it}^{*}
  =
  \rho d_{i,t-1}^{*}
  +\beta'X_{it}
  +\gamma'M_t
  +\lambda_{g(i),t}
  +\alpha_i
  +\eta_{it},
\]
where $\lambda_{g(i),t}$ is an industry- or geography-by-time effect and
$\alpha_i$ is persistent firm heterogeneity. The focal-bank observation is a
measurement equation:
\[
  \log a_{if,\mathrm{dep},t}^{\mathrm{obs}}
  =
  d_{it}^{*}
  +\log s_{if,\mathrm{dep},t}
  +\varepsilon_{it}^{D}.
\]
This equation shows the identification problem directly: observed focal-bank
deposits identify $d_{it}^{*}+\log s_{ift}$, not its two components, unless the
model has share anchors, conversion observations, account-aggregation data, or
credible restrictions on $s_{ift}$. Dynamic-panel estimators, hierarchical
state-space models, and stress-regime interactions are then estimation choices;
the statistical interface is the predictive distribution
\[
  p(D_{i,t+h}^{*}\mid \mathcal{I}_{it},\mathcal{A}_{it}),\qquad h=1,\ldots,H.
\]
Arellano-Bond-style dynamic panels are useful when many SMEs are observed for a
moderate number of periods and lagged balances interact with firm effects
\citep{arellano1991panel}.

\paragraph{Module C: discrete choice and deposit share.}
Given draws of $h_{it}$ and $D_{it}^{*}$, the choice module estimates how wallet
is allocated across the focal bank, competitors, residual unmodeled banks, and
true outside options. In a bank-level deposit model, $j$ indexes bank
relationships plus residual and outside alternatives. In a bank-product model,
$j$ indexes bank-product bundles. The interpretation of shares and diversion
must follow that choice-set definition. For alternative $j$ in market $m$,
\[
  U_{ijmt}(o)
  =
  \delta_{jmt}
  +\beta_i'A_{jmt}
  -\alpha_iP_{ijmt}(o)
  +\kappa_i'B_{ijmt}
  +\xi_{jmt}
  +\epsilon_{ijmt},
\]
so that
\[
  \Pr(c_{it}=j\mid o,\mathcal{I}_{it})
  =
  \int
  \frac{\exp(V_{ijmt}(o;\nu))}
       {\sum_{\ell\in\mathcal{J}_{imt}}\exp(V_{i\ell mt}(o;\nu))}
  \,dF_{\nu}(\nu).
\]
The focal-bank product share is
\[
  s_{ifpt}(o)
  =
  \sum_{j\in\mathcal{J}_{fpt}}
  \Pr(c_{it}=j\mid o,\mathcal{I}_{it}).
\]
BLP-style models supply the counterfactual substitution structure
\citep{berry1995automobile}; simulation methods are standard for heterogeneous
choice models \citep{train2009discrete}. Ma's framework is relevant when
$P_{ijmt}$ is individual-specific and unchosen prices must be inferred or
selection-corrected \citep{ma2024discrete}. The module must output not only
$s_{ifpt}(o)$ but also the diversion matrix
\[
  \mathcal{D}_{jk}(x)
  =
  -\frac{\partial s_k(x)/\partial P_j}
        {\partial s_j(x)/\partial P_j},
\]
where $x$ denotes the segment, market, or individual state at which diversion is
evaluated. Because deposit capture depends on where balances are diverted from,
the diversion matrix should add up over the same alternative set used in the
choice model. In the main deposit use case, it is used to distinguish balances
captured from modeled competitors, residual unmodeled banks, true outside
options, and balances shifted within the focal bank.

\paragraph{Module D: move-over amount and causal response.}
The choice model predicts substitution under specified prices and attributes;
it does not by itself identify the effect of marketing contact or
relationship-manager action. For an offer $o$, define
\[
  e_o(x)=\Pr(O_i=o\mid \mathcal{I}_{it}=x),\qquad
  m_o(x)=\mathbb{E}[V_i\mid O_i=o,\mathcal{I}_{it}=x].
\]
Under consistency, overlap, and randomized or conditionally ignorable assignment
given $\mathcal{I}_{it}$, an orthogonal score for the value difference between
offer $o$ and no offer is
\[
\begin{aligned}
  \psi_i(o)
  &=
  m_o(X_i)-m_0(X_i)
  +\frac{1\{O_i=o\}}{e_o(X_i)}\{V_i-m_o(X_i)\}\\
  &\quad
  -\frac{1\{O_i=0\}}{e_0(X_i)}\{V_i-m_0(X_i)\}.
\end{aligned}
\]
This score is a debiased estimating equation for an average value contrast, or
for a segment contrast after pre-specified conditioning. Turning it into
individualized uplift requires a separate heterogeneous-effect step, such as a
causal forest or an orthogonalized regression of the score on the decision-time
state, and that step should be validated as a policy rule rather than treated as
automatic identification of customer-level treatment effects. This connects
classical
propensity-score logic \citep{rosenbaum1983propensity}, uplift modeling
\citep{lo2002true}, causal forests \citep{wager2018causalforests}, and
orthogonal/double machine learning \citep{chernozhukov2018dml}. The module
outputs a stated estimand---average lift for a segment, conditional lift
$\widehat{\tau}_{o}^{Y}(x)$, or an individualized policy score
$\widehat{J}_{io}$---with uncertainty, overlap diagnostics, and a flag
indicating whether the estimate is experimental, quasi-experimental, or only
associational.

For product-size decisions, the outcome vector should decompose the observed
change:
\[
  Y_{iopH}
  =
  \left(
    \Delta a_{iop}^{\mathrm{move}},
    \Delta a_{iop}^{\mathrm{expand}},
    \Delta a_{iop}^{\mathrm{cann}},
    \Delta \mathrm{NPV}_{iop}
  \right).
\]
The move-over estimand is
\[
  \mu_{iop}(H)
  =
  \mathbb{E}\!\left[
  \Delta a_{iop}^{\mathrm{move}}(H)
  \mid \mathcal{I}_{it},o
  \right],
  \qquad
  \widetilde{m}_{iop}(H)=
  \frac{\mu_{iop}(H)}
       {\mathbb{E}[A_{ipt}^{-f}\mid\mathcal{I}_{it}]},
\]
where $\widetilde{m}_{iop}(H)$ is an operational reporting ratio, not generally
the same object as the structural fraction $M_{iop}(H)$. A ratio of posterior
means need not equal the posterior mean of a latent fraction under
heterogeneity. The ratio should be truncated to $[0,1]$ only as a reporting
convention and only when the denominator is identified. When data do not
separately observe migrated competitor balances, the model should report net
incremental balance, not label it move-over. The
required training data are events that identify origin or displacement:
account-transfer records, payoff of competitor debt, balance-transfer activity,
pre/post account aggregation, direct-deposit or ACH migration, or randomized
offers with follow-up account consolidation evidence.

\paragraph{Module E: policy optimization and uncertainty propagation.}
The policy layer consumes simulation draws, not point predictions:
\[
  \left\{
  D_{i,t+h}^{*(s)}, W_{ip,t+h}^{*(s)},
  A_{ip,t+h}^{-f(s)}, s_{ifp,t+h}^{(s)}(o),
  M_{iop}^{(s)}(H), \Delta a_{iop}^{\mathrm{move}(s)},
  \tau_{io}^{Y,(s)}, R_{io}^{(s)}
  \right\}_{s=1}^{S}.
\]
For each feasible offer,
\[
  \widehat{J}_{io}(H)
  =
  \frac{1}{S}\sum_{s=1}^{S}
  \left[
    V_{i,t+H}^{(s)}(o)-V_{i,t+H}^{(s)}(0)
  \right].
\]
The recommended action is
\[
  \widehat{\pi}(\mathcal{I}_{it})
  =
  \arg\max_{o\in\mathcal{O}_{it}}\widehat{J}_{io}(H)
\]
subject to chance and resource constraints such as
\[
  \Pr(R_{io}>r_{\max}\mid\mathcal{I}_{it})\le\alpha_R,\qquad
  \sum_i 1\{\widehat{\pi}(\mathcal{I}_{it})=o\}\le C_o.
\]
The required policy output is the chosen offer, expected incremental value,
credible interval, second-best action, margin of victory, and binding
constraints. If the credible interval for the margin of victory includes zero,
the system should recommend experimentation or relationship-manager review
rather than automatic targeting.

\subsection{Potential Issues and Remedies: Identification Gates}

The model should pass explicit gates before it is used for material pricing or
targeting. Each gate has a null failure mode, a diagnostic statistic, and a
permitted action if the gate fails.

\paragraph{Gate 1: latent wallet scale.}
Null failure mode: the model has learned relative customer size but not the
level of competitor wallet. Let $\mathcal{V}_{gp}$ be an external validation
set for product $p$ and segment $g$. Define relative bias and interval coverage:
\[
  B_{gp}
  =
  \frac{1}{|\mathcal{V}_{gp}|}
  \sum_{i\in\mathcal{V}_{gp}}
  \frac{\widehat{W}_{ipt}^{*}-W_{ipt}^{A}}{W_{ipt}^{A}},
  \qquad
  C_{gp}(\alpha)
  =
  \frac{1}{|\mathcal{V}_{gp}|}
  \sum_{i\in\mathcal{V}_{gp}}
  1\{W_{ipt}^{A}\in I_{ipt}^{1-\alpha}\}.
\]
The gate passes only if $|B_{gp}|\le b_{\max}$ and
$|C_{gp}(\alpha)-(1-\alpha)|\le c_{\max}$ for material segments. If it fails,
the model may be used for ranking or exploratory relationship planning, but not
for dollar wallet claims. The remedy is more anchor data, a simpler segment
model, or a sensitivity interval for $W_{ipt}^{*}$
\citep{kamakura1991latent,dempster1977em}.

\paragraph{Gate 2: missing-not-at-random competitor wallet.}
Null failure mode: missing competitor information is systematically related to
the unobserved wallet. One sensitivity model is
\[
  \operatorname{logit}\Pr(R_{ipt}^{A}=1\mid X_{it},W_{ipt}^{*})
  =
  \zeta'X_{it}+\lambda\log W_{ipt}^{*},
\]
where $R_{ipt}^{A}$ indicates anchor observation and $\lambda$ controls
selection on the latent wallet. The diagnostic is the policy range
\[
  \Delta_{\Lambda}
  =
  \sup_{\lambda\in\Lambda}\widehat{J}_{io}^{(\lambda)}
  -
  \inf_{\lambda\in\Lambda}\widehat{J}_{io}^{(\lambda)}.
\]
If the recommended action changes for plausible $\lambda$, the action is not
identified from the available data. The remedy is to collect validation anchors
or restrict the deployment to segments where the policy is stable over
$\Lambda$.

\paragraph{Gate 3: price selection.}
Null failure mode: unchosen prices are missing selectively. Let $R_{ijmt}^{P}$
indicate that the effective price for alternative $j$ is observed. A selection
correction requires variables $Z_{ijmt}^{P}$ that affect price observation or
quote generation but do not directly enter demand after conditioning on
$\mathcal{I}_{it}$:
\[
  R_{ijmt}^{P}
  =
  1\{\pi'Z_{ijmt}^{P}+v_{ijmt}>0\}.
\]
The diagnostic is whether $R_{ijmt}^{P}$ remains predictable from risk, size,
industry, or relationship variables after the proposed correction. If no
credible exclusion or randomized quote variation exists, price elasticities
should be labeled associational. Ma's selectivity correction and Heckman's
sample-selection framework give the relevant templates
\citep{heckman1979selection,ma2024discrete}.

\paragraph{Gate 4: endogenous prices and offers.}
Null failure mode: $E[C_{ijmt}\xi_{jmt}]\neq0$, so estimated price sensitivity
loads on unobserved demand or unobserved quality. With instruments $Z_{jmt}$,
the core moment is
\[
  \mathbb{E}[Z_{jmt}\xi_{jmt}(\theta)]=0.
\]
Diagnostics include first-stage strength, sign stability under alternative
instrument sets, placebo effects before rate-sheet changes, and comparison with
randomized offer cells. If the instrument set fails, the model may simulate
attribute preference but should not be used for rate optimization. Remedies are
randomized price tests, cost or rate-sheet instruments, policy-threshold
variation, or a reduced-form campaign experiment
\citep{berry1995automobile,dick2008demand}.

\paragraph{Gate 5: dynamic forecast validity.}
Null failure mode: the deposit model confuses persistence, seasonality, and
wallet growth. Because $D_{it}^{*}$ is latent for most SMEs, the gate must be
defined on an observable implication unless an anchored validation sample is
available. For horizon $h$, compare the model loss with a transparent baseline
class $\mathcal{B}$:
\[
  R_h
  =
  \frac{\sum_{(i,t)\in\mathcal{T}_{\mathrm{test}}}
    L(a_{if,\mathrm{dep},t+h}^{\mathrm{obs}},
      \widehat{D}_{i,t+h}^{*}\widehat{s}_{if,\mathrm{dep},t+h})}
  {\min_{b\in\mathcal{B}}\sum_{(i,t)\in\mathcal{T}_{\mathrm{test}}}
    L(a_{if,\mathrm{dep},t+h}^{\mathrm{obs}},
      \widehat{a}_{if,\mathrm{dep},t+h}^{\,b})}.
\]
For anchored samples with observed or credible proxy total demand
$D_{i,t+h}^{A}$, a separate latent-demand loss
$L(D_{i,t+h}^{A},\widehat{D}_{i,t+h}^{*})$ should be reported. The two losses
answer different questions and should not be pooled.
The gate requires $R_h<1$ for the horizons used in decisioning, calibrated
prediction intervals, no material residual autocorrelation, and separate
stress-window performance. If it fails, use the dynamic model only as a feature
generator and fall back to the best transparent baseline. Dynamic-panel tests
and lagged dependent variable methods are relevant here
\citep{arellano1991panel}.

\paragraph{Gate 6: substitution and cannibalization.}
Null failure mode: the choice model predicts adoption but not where balances
come from. The diversion matrix
\[
  \mathcal{D}_{jk}
  =
  -\frac{\partial s_k/\partial P_j}{\partial s_j/\partial P_j}
\]
should have economically plausible signs and satisfy adding-up over the
inside and outside alternatives. The gate also requires post-offer accounting:
\[
  \Delta a_{ifp}^{\mathrm{net}}
  =
  \Delta a_{ifp}^{\mathrm{new}}
  -\Delta a_{ifp}^{\mathrm{cannibalized}}.
\]
If diversion or cannibalization cannot be measured, the model may rank deposit
interest but should not forecast incremental captured balances. Remedies include
nested choice sets, random coefficients, explicit outside options, and deposit
offer experiments \citep{berry1995automobile,train2009discrete}.

\paragraph{Gate 7: move-over amount identification.}
Null failure mode: the model observes focal-bank growth after an offer and
mislabels it as competitor wallet moved over. Let
$Z_{iopH}^{\mathrm{mig}}$ denote migration-origin evidence, such as account
transfer records, payoff of competitor debt, balance-transfer activity,
account-aggregation before and after conversion, or direct-deposit and ACH
rerouting. The measurement equation is
\[
  Z_{iopH}^{\mathrm{mig}}
  =
  \Delta a_{iop}^{\mathrm{move}}(H)+\epsilon_{iopH}^{\mathrm{mig}},
\]
when such evidence is observed. The gate requires the reconciled decomposition
\[
  \Delta a_{ifp,t:t+H}^{\mathrm{obs}}
  =
  \widehat{\Delta a}_{iop}^{\mathrm{move}}
  +\widehat{\Delta a}_{iop}^{\mathrm{expand}}
  -\widehat{\Delta a}_{iop}^{\mathrm{cann}}
  +\widehat{u}_{iopH}
\]
to have small residuals in validation samples, economically plausible signs,
and the cap
\[
  0\leq \widehat{\Delta a}_{iop}^{\mathrm{move}}
  \leq \widehat{A}_{ipt}^{-f}
\]
except for measured new-product demand or measurement-error allowances. If this
gate fails, the bank may report net incremental balances but should not report
``dollars moved from competitors.'' Remedies are explicit migration questions in
onboarding, account-aggregation validation samples, balance-transfer and payoff
tracking, relationship-manager disposition codes, and randomized consolidation
offers.

\paragraph{Gate 8: causal incrementality.}
Null failure mode: high predicted propensity is mistaken for treatment effect.
For each offer, require overlap
\[
  \epsilon < e_o(\mathcal{I}_{it}) < 1-\epsilon
\]
on the deployment population, standardized mean differences below a chosen
threshold after weighting or matching, no placebo effect on pre-treatment
outcomes, and decision-relevant lift estimates in randomized holdouts with
calibrated uncertainty. The requirement is credible identification and
calibration, not uniformly positive lift in every subgroup. If the gate fails,
the output must be labeled propensity or association, not uplift. Randomized
holdouts are the preferred remedy; when they are unavailable, use propensity
score adjustment, causal forests, or orthogonal scores with explicit overlap
and placebo diagnostics
\citep{rosenbaum1983propensity,lo2002true,wager2018causalforests,chernozhukov2018dml}.

\paragraph{Gate 9: risk-adjusted policy validity.}
Null failure mode: the model maximizes balance while violating profitability,
liquidity, credit, conduct, or fairness constraints. The policy must satisfy
constraints such as
\[
  \mathbb{E}[\mathrm{NPV}_{io}\mid\mathcal{I}_{it}]>0,\quad
  \Pr(\mathrm{Runoff}_{io}>r_{\max}\mid\mathcal{I}_{it})\le\alpha_R,
  \quad
  |\mathrm{Calib}_g-\mathrm{Calib}_{g'}|\le\delta
\]
for monitored groups $g,g'$, where $\mathrm{Calib}_g$ is the chosen calibration
error for group $g$ on the governed target, such as predicted incremental value,
approval probability, balance lift, or adverse outcome rate. Diagnostics include
balance decay, runoff under stress, credit migration, complaint rates, subgroup
calibration, and relationship-manager override patterns. If this gate fails,
targeting should be limited to manual review or randomized learning cells until
the binding risk is understood.

\section{Deposit Pricing and Stress-Condition Demand Extension}

The recommended model in Section 5 gives the baseline architecture:
\[
  D_{ift}(o,M_t)=D_{it}^{*}(M_t)\,s_{ift}(o,M_t).
\]
This section extends that architecture for deposit pricing and stress
conditions. The extension is not a separate business question. It answers the
same wallet-capture question under a richer environment: given market condition
$M_t$, stress state $R_t$, and deposit action $o$, how much total SME deposit
demand exists, what share can the focal bank capture, what price is required,
and how stable is the captured balance?
The stress-pricing problem is therefore a joint problem of \emph{quantity},
\emph{share}, \emph{price}, and \emph{duration}. A high offered rate is
attractive only if the resulting balances are incremental, sufficiently stable,
and valuable after funding, liquidity, cannibalization, and relationship costs.

\subsection{Pricing Objective and Decision Estimands}

Let $r_{io}$ denote the effective deposit rate, earnings credit, fee waiver, or
relationship-pricing term embedded in action $o$. A pricing decision should
maximize incremental relationship value, not gross balance:
\[
\begin{aligned}
  J_{io}^{\mathrm{price}}(H)
  =
  \mathbb{E}\Big[
  &\sum_{h=1}^{H}
    \delta^h
    \{v_{t+h}^{F}D_{if,t+h}(o,R_{t+h})
    - c^{r}_{io,t+h}D_{if,t+h}(o,R_{t+h}) \\
  &\quad
    + \Pi_{i,t+h}^{\mathrm{rel}}(o)
    - \mathrm{EL}_{i,t+h}(o)
    - \mathrm{LC}_{i,t+h}(o)\}
    - C_{io}
  \mid \mathcal{I}_{it}\Big],
\end{aligned}
\]
where $v_{t+h}^{F}$ is the funding or liquidity value of deposits, $c^{r}_{io}$
is the rate or pricing cost, $\Pi^{\mathrm{rel}}$ is non-deposit relationship
profit, $\mathrm{EL}$ is expected credit loss, $\mathrm{LC}$ is liquidity or
runoff cost, and $C_{io}$ is campaign or onboarding cost. The decision rule is
\[
  o_i^{*}
  =
  \arg\max_{o\in\mathcal{O}_{it}}
  J_{io}^{\mathrm{price}}(H)
\]
subject to conduct, fair-lending, liquidity, concentration, relationship-manager
capacity, and profitability constraints. This prevents the pricing model from
recommending expensive balances that are transient, risky, or cannibalized from
existing focal-bank deposits.

The target should be an \emph{incremental} pricing estimand. Let
\[
  \Delta D_{ifo,t+h}^{r}
  =
  D_{if,t+h}^{r}(o)-D_{if,t+h}^{r}(0)
\]
denote incremental focal-bank deposits in scenario $r$. The object to be priced
is not this total alone, but the decomposition
\[
  \Delta D_{ifo,t+h}^{r}
  =
  \Delta D_{ifo,t+h}^{\mathrm{move},r}
  +\Delta D_{ifo,t+h}^{\mathrm{expand},r}
  -\Delta D_{ifo,t+h}^{\mathrm{cann},r}
  +\epsilon_{io,t+h}^{r}.
\]
The first term is competitor wallet moved to the focal bank, the second is
growth in the SME's total liquidity demand, and the third is cannibalization of
existing focal-bank balances or organic inflows. Stress pricing should report
all three terms. Otherwise, the bank may pay a premium rate for balances it
already would have retained or for temporary balances that leave as soon as
market conditions normalize.

\subsection{Economic Pricing Logic and Marginal Value}

A useful way to discipline the pricing decision is to write the period
contribution margin as
\[
  \pi_{io,t+h}^{r}
  =
  v_{t+h}^{F,r}D_{if,t+h}^{r}(o)
  -r_{io,t+h}D_{if,t+h}^{r}(o)
  -\Phi_{io,t+h}^{r}(D_{if,t+h}^{r},o)
  +\Pi_{io,t+h}^{\mathrm{rel},r},
\]
where $\Phi$ includes servicing cost, incentive cost, liquidity charge,
operational cost, and runoff or concentration penalties. Holding the feasible
offer class fixed, an interior rate choice must satisfy
\[
\begin{aligned}
  0
  =
  \mathbb{E}\!\Bigg[
  \sum_{h=1}^{H}\delta^h
  \Big\{
    &(v_{t+h}^{F,r}-r_{io,t+h})
    \frac{\partial D_{if,t+h}^{r}}{\partial r_{io}}
    -D_{if,t+h}^{r} \\
    &-\frac{\partial \Phi_{io,t+h}^{r}}{\partial r_{io}}
    +\frac{\partial \Pi_{io,t+h}^{\mathrm{rel},r}}{\partial r_{io}}
  \Big\}
  \mid\mathcal{I}_{it}\Bigg].
\end{aligned}
\]
This condition is useful for management because it separates the \emph{marginal
balance effect} from the \emph{infra-marginal cost}: raising the rate may
attract or retain deposits, but the higher rate is usually paid on the full
eligible balance, not only on the marginal dollars. In stress conditions this
tradeoff becomes sharper because $v^{F,r}$, runoff cost, and competitor price
responses can change at the same time. The condition also shows why a simple
deposit-beta forecast is insufficient: the pricing decision needs the dollar
response, the margin on those dollars, and the stability of those dollars.

\subsection{Integrating Pricing with Discrete Choice}

In the Section 5 discrete-choice model, separate cost-like price terms from
quality and safety attributes. Let the effective deposit-price vector be
\[
  C_{ijmt}(o)
  =
  \big(r_{ijmt}(o), f_{ijmt}(o), \mathrm{minbal}_{ijmt}(o),
  \mathrm{feecredit}_{ijmt}(o)\big),
\]
where $r$ is the deposit rate or earnings credit, $f$ are fees, $\mathrm{minbal}$
is the minimum-balance requirement, and $\mathrm{feecredit}$ captures explicit
or implicit relationship pricing. Let
\[
  Q_{jmt}(o)
  =
  \big(\mathrm{safety}_{jmt}(o),\mathrm{service}_{jmt}(o),
  \mathrm{digital}_{jmt}(o),\mathrm{treasury}_{jmt}(o)\big)
\]
collect quality attributes. The safety component includes deposit insurance
structure, collateralization, sweep arrangements, operational continuity,
perceived bank soundness, and relationship-manager credibility.
Classic bank-run theory makes the safety and liquidity dimension central
\citep{diamond1983bankruns}; depositor-level evidence also shows that
relationships and depositor networks matter during runs \citep{iyer2012bankruns}.
The utility for deposit relationship $j$ is
\[
  U_{ijmt}(o)
  =
  \delta_{jmt}
  +\beta_i'A_{jmt}
  +\zeta_i'Q_{jmt}(o)
  -\alpha_i'C_{ijmt}(o)
  +\kappa_i'B_{ijmt}
  +\xi_{jmt}
  +\epsilon_{ijmt}.
\]
The focal-bank share response is then
\[
  s_{ift}(o)
  =
  \sum_{j\in\mathcal{J}_{fmt}}
  \Pr(c_{it}=j\mid C_{ijmt}(o),Q_{jmt}(o),A_{jmt},B_{ijmt},\mathcal{I}_{it}).
\]
The balance response to pricing is not the same as the choice probability. It
is the product of total demand and share:
\[
  \frac{\partial D_{ift}}{\partial r_{io}}
  =
  \frac{\partial D_{it}^{*}}{\partial r_{io}}s_{ift}
  +
  D_{it}^{*}\frac{\partial s_{ift}}{\partial r_{io}}.
\]
For ordinary operating balances, the first term may be small if the price
mainly reallocates balances across banks. During stress, the first term can be
material if the pricing or safety signal affects the SME's desired liquidity.
This equation is the reason stress pricing must be combined with the main
dynamic-demand and discrete-choice model.

The same equation also clarifies how to distinguish three common pricing
stories. If $\partial D_{it}^{*}/\partial r_{io}\approx 0$ and
$\partial s_{ift}/\partial r_{io}>0$, pricing mainly reallocates existing SME
liquidity from competitors. If both terms are positive, the offer both attracts
share and induces larger desired balances, perhaps through perceived liquidity
support or safety. If the share term is positive but the post-offer survival of
balances is low, the bank has bought rate-sensitive hot money rather than a
durable operating relationship.

\subsection{Heterogeneous Rate Sensitivity and Deposit Beta}

A reduced-form pricing equation useful for implementation is
\[
  \log D_{ifmt}
  =
  \alpha_i+\delta_m+\lambda_t
  +\beta_{r,i}(R_t) r_{ifmt}
  +\beta_c C_{fmt}
  +\beta_q Q_{fmt}
  +\beta_b B_{ift}
  +u_{ifmt},
\]
where $m$ is a local market, $R_t$ is the regime, $C_{fmt}$ is local competition,
$Q_{fmt}$ captures branch, digital, and service quality, and $B_{ift}$ captures
relationship depth and switching frictions. The rate coefficient should vary by
customer and regime:
\[
  \beta_{r,i}(R_t)
  =
  \beta_0
  +\beta_x'X_i
  +\beta_m'M_m
  +\beta_R R_t
  +\beta_{xR}'(X_i\otimes R_t)
  +\zeta_i.
\]
This specification lets the model represent SMEs whose operating balances are
rate insensitive in normal periods but become rate sensitive during liquidity
competition, as well as SMEs that prioritize safety or operating continuity
over rate. Deposit-competition and monetary-policy papers motivate this focus
on market power, rate pass-through, and deposit betas
\citep{hannan1991rigidity,neumark1992market,egan2017deposit,drechsler2017deposits};
structural banking demand
papers motivate the explicit role of prices, attributes, and local competition
\citep{dick2008demand,adams2007who}.

For implementation, it is important to separate the \emph{deposit beta} from
the \emph{customer balance elasticity}. A deposit beta is a pass-through object,
for example
\[
  b_{if}^{r}
  =
  \frac{\partial r_{ifmt}}{\partial r_{t}^{\mathrm{mkt}}}
  \quad\text{or}\quad
  \frac{\Delta r_{ifmt}}{\Delta r_{t}^{\mathrm{mkt}}},
\]
where $r_t^{\mathrm{mkt}}$ is a market policy or wholesale-rate benchmark. The
balance elasticity is
\[
  \varepsilon_{if}^{r}
  =
  \frac{\partial \log D_{ifmt}^{r}}{\partial r_{ifmt}}.
\]
Stress pricing needs both. A high beta can be a defensive pass-through policy
with little balance response, while a high balance elasticity can imply a large
response to small pricing differences. The recommended model should therefore
estimate a response surface
\[
  \varepsilon_{if}^{r}
  =
  e(X_i,M_m,R_t,\mathrm{Safety}_{fm,t},
  \mathrm{SwitchCost}_{if,t},\mathrm{Uninsured}_{i,t}),
\]
where $\mathrm{Uninsured}_{i,t}$ measures the portion of balances exposed to
deposit-insurance limits or uninsured operational risk when the bank has such
data or can approximate it. Monetary-tightening stress and uninsured depositor
run evidence make this exposure especially relevant
\citep{jiang2023monetary}.

\subsection{Branch Market Share as a Noisy Anchor}

The bank may observe branch-level deposit market share at granular geography,
but such shares usually mix consumer, SME, commercial, and wholesale deposits.
They should therefore anchor local competition and broad share, not be treated
as direct SME labels. Let $\mathcal{G}$ denote segments such as consumer, SME,
commercial, and wholesale. A measurement equation is
\[
  S_{fmt}^{\mathrm{branch}}
  =
  \sum_{g\in\mathcal{G}}\omega_{gmt}S_{fgmt}
  +e_{fmt},
\]
where $S_{fgmt}$ is segment-specific focal-bank share and $\omega_{gmt}$ is the
local deposit mix. If $\omega_{gmt}$ is not directly observed, the model should
estimate it using branch location, account mix, relationship-manager portfolio,
public market data, and internal segment tags. The deployment output should
distinguish:
\[
  \Pr(S_{f,\mathrm{SME},m,t}\mid S_{fmt}^{\mathrm{branch}},\mathcal{I}_{mt})
  \quad\text{from}\quad
  S_{fmt}^{\mathrm{branch}}.
\]
Confusing these two objects would overstate how much branch-level share
identifies SME wallet share.

\subsection{Stress Regimes and Deposit Demand}

Stress conditions change both total deposit demand and focal-bank share. Let
\[
  R_t\in
  \{\mathrm{normal},\mathrm{rate\ shock},\mathrm{credit\ stress},
  \mathrm{liquidity\ stress},\mathrm{flight\ to\ safety}\}
\]
denote the regime. The total-demand layer can be written as
\[
\begin{aligned}
  \log D_{i,t+h}^{*}
  &=
  \rho \log D_{it}^{*}
  +\gamma'X_{i,t+h}
  +\lambda_{g(i),t+h}
  +\theta_R'R_{t+h} \\
  &\quad
  +\theta_E'\mathrm{Exposure}_{i,t+h}
  +\theta_L'\mathrm{Liquidity}_{i,t+h}
  +\eta_i+\nu_{i,t+h}.
\end{aligned}
\]
The share layer should then allow the same regime to alter price sensitivity,
safety preference, and switching frictions. It is cleaner to write the stress
utility directly from primitives rather than adding stress terms to a baseline
utility that may already contain price:
\[
\begin{aligned}
  U_{ijm,t+h}^{R}(o)
  &=
  \delta_{jm,t+h}^{R}
  +\beta_i^{R\prime}A_{jm,t+h}
  +\zeta_i^{R\prime}Q_{jm,t+h}(o)
  -\alpha_i^{R\prime}C_{ijm,t+h}(o) \\
  &\quad
  +\kappa_i^{R\prime}B_{ijm,t+h}
  +\chi_C'(R_{t+h}\otimes C_{ijm,t+h}(o))\\
  &\quad
  +\chi_Q'(R_{t+h}\otimes Q_{jm,t+h}(o))
  +\chi_S'(R_{t+h}\otimes \mathrm{SwitchCost}_{ijm,t+h})
  +\xi_{jm,t+h}^{R}
  +\epsilon_{ijm,t+h}.
\end{aligned}
\]
Any regime term common to all alternatives is normalized into the outside
option or omitted, because it cancels in a standard logit share. A regime
matters for share only through alternative-specific attributes, interactions,
or an explicit outside-option normalization.
During stress, the sign of deposit change is not universal. Some SMEs may build
balances through precautionary saving, policy inflows, or credit-line draws;
others may burn cash as revenue falls. Evidence from COVID-period corporate
studies documents liquidity draws and a ``dash for cash''
\citep{acharya2020dash,li2020banks}; broader theory links bank lending
commitments and deposit-taking \citep{kashyap2002liquidity}; and credit-line
research treats lines as liquidity insurance
\citep{sufi2009creditlines,acharya2014creditlines}. For SMEs, the model should
estimate this heterogeneity by industry, cash-flow exposure, credit access,
relationship depth, and geography rather than imposing a single stress
multiplier.

For stress pricing, the regime variable should not be only a historical label.
It should be part of a scenario generator. One implementation is a transition
model
\[
  \Pr(R_{t+1}=k\mid R_t,\mathcal{Z}_t)
  =
  \frac{\exp(a_{kR_t}+b_k'\mathcal{Z}_t)}
       {\sum_{\ell}\exp(a_{\ell R_t}+b_\ell'\mathcal{Z}_t)},
\]
where $\mathcal{Z}_t$ includes market rates, local unemployment, industry
stress, credit spreads, deposit outflow indicators, competitor rate gaps, and
bank-specific safety indicators. A management stress scenario can override this
transition model by conditioning on a path
$\{R_{t+1},\ldots,R_{t+H}\}$ and corresponding paths for rates, credit
conditions, and competitor actions. This keeps the statistical model and the
management scenario internally consistent.

\subsection{Runoff, Duration, and Stable Balance Value}

Stress-condition pricing is incomplete without a runoff model. Let
$T_{io}^{r}$ be the time until the incremental balance generated by offer $o$
runs off in scenario $r$. Define the survival curve
\[
  S_{io}^{r}(h)
  =
  \Pr(T_{io}^{r}>h\mid
  \mathcal{I}_{it},o,R_{t:t+h}=r)
\]
and the stable-balance contribution
\[
  \mathrm{SBD}_{io}^{r}(H)
  =
  \sum_{h=1}^{H}\delta^h
  S_{io}^{r}(h)\,
  \Delta D_{ifo,0}^{\mathrm{move},r}.
\]
Here $\Delta D_{ifo,0}^{\mathrm{move},r}$ is the originated moved balance
before runoff. If the model instead forecasts a realized post-runoff path
$\Delta D_{ifo,t+h}^{\mathrm{move},r}$, then the stable-balance contribution
should be $\sum_h\delta^h\Delta D_{ifo,t+h}^{\mathrm{move},r}$ without another
multiplication by $S_{io}^{r}(h)$. Mixing these two conventions double counts
runoff.
The model should estimate separate survival curves for moved competitor
balances, expanded liquidity demand, and cannibalized balances:
\[
  S_{io}^{k,r}(h),\qquad
  k\in\{\mathrm{move},\mathrm{expand},\mathrm{cann}\}.
\]
This distinction matters because a rate promotion may generate large first-week
or first-month inflows but weak stable-balance value. A practical hazard model
is
\[
\begin{aligned}
  &\Pr(T_{io}^{k,r}=h\mid T_{io}^{k,r}\ge h,\mathcal{I}_{it}) \\
  &\quad =
  \Lambda\!\left(
    a_{krh}
    +b_{kr}'X_i
    +c_{kr}'\mathrm{Offer}_{io}
    +d_{kr}'\mathrm{CompetitorGap}_{mt+h}
    +e_{kr}'\mathrm{Relationship}_{if,t+h}
  \right),
\end{aligned}
\]
where $\Lambda$ is a logit or complementary-log-log link. The pricing objective
should value balances by expected survival, not merely by day-one inflow.
Thus a serious stress-pricing report should show both
$\widehat{\Delta D}_{ifo}^{r}(H)$ and $\widehat{\mathrm{SBD}}_{io}^{r}(H)$.

\subsection{Stress Offer Design: Rate, Safety, and Relationship Terms}

In stress conditions, the bank should not treat ``price'' as only the posted
rate. The action vector should be
\[
  o
  =
  (r,\mathrm{fee},\mathrm{earnings\ credit},
  \mathrm{sweep},\mathrm{collateral},\mathrm{operating\ support},
  \mathrm{relationship\ terms}).
\]
For some SMEs, especially operating-account customers with payroll, payables,
and merchant settlement needs, reliability, funds availability, credit-line
access, and deposit-insurance or sweep structure can dominate a small rate
difference. For other SMEs with excess cash and low switching costs, the rate
gap may dominate. The model should therefore estimate substitution between
rate and safety or relationship terms, not only the level effect of rate.
A useful willingness-to-pay contrast is
\[
  \mathrm{WTP}_{i}^{\mathrm{safety},r}
  =
  \frac{
    \partial U_{ijmt}^{R}/\partial Q_{jmt}^{\mathrm{safety}}}
       {\partial U_{ijmt}^{R}/\partial r_{ijmt}},
\]
interpreted as the rate concession or premium that holds utility fixed for a
one-unit safety improvement, with the sign determined by the rate convention.
If higher $r$ is better for depositors, the denominator is positive; if $r$ is
entered as a bank cost, the sign should be reversed. The scale of the safety
measure must be reported for the WTP number to be meaningful.
This connects the stress extension back to the discrete-choice model and to
Ma's individual-specific financial-product pricing framework
\citep{ma2024discrete}.

\subsection{Scenario Outputs for Management}

For each stress scenario $r$ and offer $o$, the model should produce a
decomposition:
\[
\begin{aligned}
  \widehat{\Delta D}_{ifo}^{r}(H)
  &=
  \widehat{D}_{i,t+H}^{*,r}(0)\widehat{\Delta s}_{ifo}^{r}(H)
  +\widehat{s}_{if,t+H}^{r}(0)\widehat{\Delta D}_{io}^{*,r}(H) \\
  &\quad
  +\widehat{\Delta s}_{ifo}^{r}(H)
   \widehat{\Delta D}_{io}^{*,r}(H),
\end{aligned}
\]
where the first term is share capture, the second is total-demand expansion,
and the third is interaction. Management should see these as separate numbers,
not only a single incremental-balance forecast. The pricing package should also
report:
\[
  \{\widehat{\Delta D}^{r},\widehat{\mathrm{SBD}}^{r},
  \widehat{\mathrm{NPV}}^{r},
  \widehat{\mathrm{Runoff}}^{r},
  \widehat{\mathrm{Cannibalization}}^{r},
  \widehat{\mathrm{LiquidityValue}}^{r},
  \widehat{\mathrm{ConfidenceInterval}}^{r}\}.
\]
This output answers the practical question: whether a higher offered rate is
buying durable SME operating deposits, temporary hot money, balances that would
have arrived without pricing, or balances cannibalized from existing focal-bank
relationships.

The scenario package should also report break-even pricing. Because demand,
runoff, cannibalization, and relationship spillovers can make value nonmonotone
in the offered rate, the primary object should be the acceptable-rate set:
\[
  \mathcal{R}_{io}^{\mathrm{BE},r}
  =
  \{r\in\mathcal{R}_{io}:
  \widehat{J}_{io}^{\mathrm{price},r}(H)\ge 0\}.
\]
Only when this set is an interval, or when the model explicitly assumes local
monotonicity over the feasible offer range, should the report summarize it by
$r_{io}^{\mathrm{BE},r}=\sup \mathcal{R}_{io}^{\mathrm{BE},r}$. If the proposed
offer rate lies outside $\mathcal{R}_{io}^{\mathrm{BE},r}$, the model should
explain which assumption makes the action still acceptable, such as
relationship profit, strategic retention, liquidity value, or a management
override. If no such assumption is documented, the offer is not economically
justified even if it is expected to grow balances.

\subsection{Identification and Validation Requirements}

Stress pricing requires stronger evidence than normal-period ranking. The key
identification threats are endogenous pricing, endogenous relationship-manager
attention, unobserved competitor rates, segment mixing in branch market share,
and regime-specific selection. A serious implementation should use at least one
of the following sources of identifying variation:

\begin{enumerate}[leftmargin=*]
  \item randomized or quasi-random deposit-rate and fee-waiver tests;
  \item rate-sheet changes or policy thresholds that create exogenous price
  variation;
  \item competitor-rate shocks or local market shocks used as structural demand
  shifters, or as instruments only when an explicit exclusion restriction is
  credible after conditioning on local demand and quality;
  \item branch or relationship-manager capacity variation that affects contact
  but not latent deposit demand after controls;
  \item pre/post conversion evidence that identifies gross migrated balances,
  post-stress runoff, and cannibalization.
\end{enumerate}

Validation should include normal and stress windows separately. Required gates
are:
\[
  \mathrm{Bias}_{gr}
  =
  \frac{1}{|V_{gr}|}\sum_{i\in V_{gr}}
  \frac{\widehat{D}_{i}^{r}-D_i^{r}}{D_i^{r}},
  \qquad
  \mathrm{Coverage}_{gr}(\alpha)
  =
  \frac{1}{|V_{gr}|}\sum_{i\in V_{gr}}
  1\{D_i^{r}\in I_i^{r,1-\alpha}\},
\]
for segment $g$ and regime $r$. The pricing model should not be used for
material stress pricing unless dollar calibration, interval coverage,
rate-elasticity signs, runoff forecasts, and cannibalization decomposition pass
in the relevant segment. If those gates fail, the model can still support
learning-cell design and management overlays, but not automated stress pricing.

For governance, stress pricing should have explicit promotion thresholds:
\[
\begin{gathered}
  |\mathrm{Bias}_{gr}|\le b_{\max},\qquad
  |\mathrm{Coverage}_{gr}(\alpha)-(1-\alpha)|\le c_{\max},\\
  \operatorname{sign}(\widehat{\varepsilon}_{if}^{r})
    \ \text{stable across accepted specifications},\qquad
  \Pr(\widehat{\mathrm{NPV}}_{io}^{r}<0)\le q_{\max},\\
  \left|
  \widehat{\mathrm{Runoff}}_{gr}^{\mathrm{stress}}
  -
  \mathrm{Runoff}_{gr}^{\mathrm{realized}}
  \right|\le u_{\max}.
\end{gathered}
\]
The last line is especially important: a model that captures first-month
inflows but misses post-stress decay can overstate the value of stress pricing.
When historical stress observations are sparse, validation should use
backtests across rate-hiking cycles, local market shocks, industry downturns,
competitor rate changes, and deliberately randomized learning cells. The result
should be labeled by evidence grade: randomized, quasi-experimental,
historically backtested, scenario-only, or management overlay.

\section{Estimation Problems Specific to the Recommended Model}

The proposed architecture is useful only if the bank estimates the right
objects. A large premium bank faces a particularly important identification
problem: its observed SME customers are not a random sample of the SME market.
They may be larger, more complex, more urban, more digitally active, more
relationship-oriented, more multi-product, or more sensitive to brand and
operational reliability than SMEs served by regional and community banks. If the
model is fit only on focal-bank customers, it can learn the behavior of SMEs
that already selected a premium bank, not the behavior of competitor-bank
customers the bank wants to win.

This section states the estimation problems in the notation of Sections 5 and
6. The main point is that selection, unobserved quality, and endogenous prices
enter every component of
\[
  D_{ift}(o,M_t,R_t)
  =
  D_{it}^{*}(M_t,R_t)\,s_{ift}(o,M_t,R_t),
\]
and therefore directly affect estimated competitor wallet, pricing response,
move-over, and risk-adjusted value.

The closest literatures each solve part of the problem, but none solves the
full SME wallet-acquisition problem alone. BLP-style banking demand models use
market shares, product attributes, and instruments to handle unobserved
product-market quality \citep{berry1995automobile,dick2008demand,nevo2001cereal}.
Ma extends this logic to financial products with individual-specific prices and
selectively observed chosen-product prices \citep{ma2024discrete}. Sample
selection and external-validity literatures clarify why estimates from one
selected population need not transport to another
\citep{heckman1979selection,hotz2005predicting,manski1990bounds}. Causal
response methods separate propensity from incremental treatment effects
\citep{rosenbaum1983propensity,wager2018causalforests,chernozhukov2018dml}.
Our setting must combine all of these: a selected premium-bank panel, latent
competitor wallets, negotiated SME prices, branch-level aggregate anchors, and
campaign interventions.

\begin{table}[htbp]
\centering
\footnotesize
\begin{tabular}{@{}
  >{\raggedright\arraybackslash}p{0.19\linewidth}
  >{\raggedright\arraybackslash}p{0.25\linewidth}
  >{\raggedright\arraybackslash}p{0.25\linewidth}
  >{\raggedright\arraybackslash}p{0.20\linewidth}
@{}}
\toprule
Model component & What the literature typically does & Why it is insufficient here & Required adaptation \\
\midrule
Share model $s_{ift}$ & BLP/micro-BLP uses market shares, product attributes,
and instruments for prices & SME bank quality and relationship value are
partly unobserved and premium-bank clients are selected & Product-market
quality controls, instruments, and target-population diagnostics \\
Price vector $C_{ijmt}$ & Financial-product demand models allow individual-
specific prices and selected observed prices & SME prices include negotiated
rates, fee waivers, treasury terms, and rejected quotes & Quote ledger,
price-imputation validation, and selectivity correction \\
Wallet $W_{ipt}^{*}$ & Cross-sell and latent-factor models infer missing
product ownership or amount & Competitor dollar balances are latent and branch
shares mix segments & External anchors, segment-mix measurement equations, and
uncertainty intervals \\
Dynamic demand $D_{it}^{*}$ & Dynamic panels forecast balances from lags and
market conditions & Premium-bank customers may have different growth,
seasonality, and stress behavior & Stratified out-of-time tests and separation
of demand growth from share movement \\
Move-over and policy $J_{io}$ & Uplift and causal models estimate treatment
effects under assignment assumptions & RM attention, pricing exceptions, and
campaigns are targeted & Randomized learning cells, overlap diagnostics, and
orthogonal scores \\
\bottomrule
\end{tabular}
\caption{How existing estimation approaches map into the recommended SME wallet model.}
\label{tab:estimation_literature_map}
\end{table}

\subsection{Premium-Bank Selection and Transportability}

Let $B_i=f$ denote that SME $i$ is observed as a focal-bank customer. The bank
usually estimates on
\[
  p(Y_i,X_i\mid B_i=f),
\]
but deployment requires predictions for prospects currently held by competitors:
\[
  p(Y_i(o),X_i\mid B_i\neq f).
\]
The selection problem is that
\[
  p(h_{it}\mid X_i,B_i=f)
  \neq
  p(h_{it}\mid X_i,B_i\neq f),
\]
where $h_{it}$ is the latent state containing liquidity need, price sensitivity,
switching cost, service preference, and relationship depth. A Chase-like bank's
customers may have systematically higher operating complexity and stronger
preference for scale, branch/digital reliability, treasury integration, and
brand safety. A model trained only on these customers may overstate the
competitor-bank prospect's willingness to move for the same offer, or understate
the rate premium needed to dislodge a relationship at a regional bank.

The remedy is not to ``adjust for bank size'' with one control variable. The
model needs a target-population correction. Let $T_i=1$ denote the target
population for a campaign or prospect universe, and let $S_i=1$ denote
membership in the estimation sample. A transparent first correction is
importance weighting:
\[
  w_i
  =
  \frac{\Pr(T_i=1\mid X_i)}
       {\Pr(S_i=1\mid X_i)}.
\]
The weighted estimating equations for any module parameter $\theta$ become
\[
  \sum_{i,t} w_i\,\psi(Y_{it},X_{it};\theta)=0.
\]
This is the analogue of reweighting or transporting estimates from one sample
to another in program-evaluation settings \citep{hotz2005predicting}. It helps
only for selection on observed variables. If unobserved service preference or
switching cost differs between premium-bank customers and regional-bank
customers, the problem becomes partially identified; the honest output is a
sensitivity range or bound, not a single point estimate \citep{manski1990bounds}.
Therefore the bank should also use external anchor samples: consented account
aggregation, onboarding statements, win/loss files, branch-market aggregates,
market-cell benchmarks, and conversion events. The target-population diagnostic
is segment calibration:
\[
  \mathrm{Bias}_{gp}^{T}
  =
  \frac{1}{|V_{gp}^{T}|}
  \sum_{i\in V_{gp}^{T}}
  \frac{\widehat{W}_{ipt}^{*}-W_{ipt}^{A}}{W_{ipt}^{A}},
\]
computed on validation samples that resemble the prospect population, not only
current focal-bank customers. If this diagnostic fails, the model should be
reported as a within-bank relationship-deepening model, not as a competitor
wallet-acquisition model.

\subsection{Unobserved Bank Quality in the Share Model}

The discrete-choice layer contains unobserved bank-market quality:
\[
  U_{ijmt}
  =
  \delta_{jmt}
  +\beta_i'A_{jmt}
  +\zeta_i'Q_{jmt}
  -\alpha_i'C_{ijmt}
  +\kappa_i'B_{ijmt}
  +\xi_{jmt}
  +\epsilon_{ijmt}.
\]
For a premium bank, $\xi_{jmt}$ is not a nuisance term. It may include brand
trust, perceived safety, treasury-service reliability, digital uptime,
relationship-manager quality, geographic convenience for owners and employees,
and the value of being with a systemically important bank. These attributes are
partly observed by customers and relationship managers but not fully observed
by the econometrician. If cost-like prices or offers are correlated with
$\xi_{jmt}$,
then
\[
  \mathbb{E}[C_{ijmt}\xi_{jmt}]\neq0,
\]
and the estimated price coefficient confounds price sensitivity with unobserved
quality. This is the standard differentiated-product demand issue in BLP-style
models \citep{berry1995automobile,nevo2001cereal}, and it is especially
important in banking because service and safety quality are central parts of
the product.

The model should therefore estimate product-market fixed effects or mean
utilities $\delta_{jmt}$, use instruments or control functions for endogenous
price, and include explicit quality proxies where possible. Candidate quality
controls include branch density, digital-service incidents, service response
times, relationship-manager coverage, treasury implementation time, complaint
rates, local deposit outflow stress, bank credit ratings, uninsured-deposit
exposure, and customer-service capacity. Candidate instruments include
rate-sheet rules, funding-cost shifters, policy thresholds, centralized pricing
changes, and cost interactions with product attributes. Competitor-rate shocks
should usually enter as demand-environment variables or competitor-response
scenarios; they are instruments only with a documented exclusion restriction.
The identifying moment is
\[
  \mathbb{E}[Z_{jmt}\xi_{jmt}(\theta)]=0.
\]
This is the same role played by instruments in differentiated-product demand
estimation: they separate price variation caused by supply or cost shifters
from price variation caused by unobserved demand quality
\citep{berry1995automobile,nevo2001cereal}. In banking applications such as
\citet{dick2008demand} and \citet{egan2017deposit}, local competition, product
attributes, and banking-market structure provide part of that identifying
variation. For SME wallet capture, the additional difficulty is that the
premium-bank brand and relationship quality are themselves central attributes,
so the quality controls must be richer than in a simple posted-rate deposit
model.
When that moment is not credible, the output should be labeled as descriptive
preference or association. It should not be used for rate optimization.
Control-function approaches are another way to handle endogenous prices in
choice models \citep{petrin2010control,wooldridge2015control}.

\subsection{Individual-Specific Prices and Selectively Observed Quotes}

For SME banking, the effective cost-like price vector
\[
  C_{ijmt}(o)
  =
  (r,f,\mathrm{minbal},\mathrm{feecredit})
\]
is often negotiated; safety and service belong in the quality vector
$Q_{jmt}(o)$. The econometrician may observe booked prices for chosen products
and perhaps some rejected quotes, but not the full menu of prices each SME
would have faced at all competitors. This is exactly why the Ma framework is
relevant: in financial products, customers can face individual-specific prices,
and using a common market price can bias estimated price sensitivity when those
prices are correlated with demographics, risk, or income
\citep{ma2024discrete}.

For the recommended model, the first-stage pricing equation should be written
for every feasible bank-product alternative:
\[
  C_{ijmt}^{k}
  =
  q_k(Z_{ijmt}^{P},X_i,A_{jmt},M_t)+v_{ijmt}^{k},
  \qquad
  k\in\{r,f,\mathrm{minbal},\mathrm{feecredit}\}.
\]
The problem is that $C_{ijmt}^{k}$ is often observed only when a quote is
generated or a product is chosen. Let $R_{ijmt}^{P}=1$ indicate observed price.
Then
\[
  R_{ijmt}^{P}
  =
  1\{\pi'Z_{ijmt}^{R}+\nu_{ijmt}>0\}.
\]
If quote observation depends on latent demand or bargaining ability,
$\mathbb{E}[v_{ijmt}^{k}\mid R_{ijmt}^{P}=1]\neq0$. Ma's estimator treats this
as a selectivity problem distinct from the classic Heckman market-entry problem:
the customer is in the market, but the econometrician observes only the chosen
product's individualized price \citep{heckman1979selection,ma2024discrete}. In
our SME setting, the same issue appears when we observe the booked treasury or
deposit terms but not the competitor terms and rejected focal-bank quotes. The
remedy is to maintain a quote ledger that records offered, rejected, expired,
and exception-approved terms, and to use selectivity correction or control
functions when unchosen prices must be imputed. For governance, the
price-imputation layer should have its own out-of-sample loss:
\[
  L_P
  =
  \frac{1}{|Q|}
  \sum_{(i,j,m,t)\in Q}
  \|C_{ijmt}^{\mathrm{obs}}-\widehat{C}_{ijmt}\|^2,
\]
where $Q$ is a held-out quote sample. Poor price prediction should veto
counterfactual pricing even if the final choice model has good in-sample fit.

\subsection{Latent Competitor Wallet and Segment Mixing}

The model uses latent competitor wallet
\[
  A_{ipt}^{-f}=W_{ipt}^{*}-a_{ifpt}^{\mathrm{obs}}.
\]
This object is not directly observed for most prospects. Public and branch
market-share data help anchor local totals, but they mix consumer, SME,
commercial, and wholesale balances. Thus the measurement equation is
\[
  S_{fmt}^{\mathrm{branch}}
  =
  \sum_{g\in\mathcal{G}}\omega_{gmt}S_{fgmt}+e_{fmt},
\]
not $S_{f,\mathrm{SME},m,t}=S_{fmt}^{\mathrm{branch}}$. If the segment weights
$\omega_{gmt}$ are misspecified, the model can incorrectly infer that a branch
market is weak in SME share when the weakness is actually consumer deposits, or
vice versa.

The remedy is a hierarchical measurement model with sensitivity analysis over
segment-mix assumptions:
\[
  p(W^{*},S,\omega\mid
  a^{\mathrm{obs}},S^{\mathrm{branch}},A^{\mathrm{anchor}},X).
\]
This is closest in spirit to latent-trait and mixed-data factor models in the
cross-selling literature \citep{kamakura1991latent,kamakura2003database}, but
the SME wallet version has a harder level-identification problem because the
missing variable is a dollar competitor balance, not only an unobserved product
affinity. EM or Bayesian data augmentation can estimate the latent structure
\citep{dempster1977em}, but the level of $W_{ipt}^{*}$ still requires anchors.
The output should include posterior or bootstrap intervals for
$A_{ipt}^{-f}$, not only a point estimate. Management reports should distinguish
three evidence grades: directly anchored competitor wallet, market-anchored
competitor wallet, and model-imputed competitor wallet. Only the first two
should be used for dollar claims in high-stakes targeting.

\subsection{Dynamic Demand Versus Share Movement}

Observed deposit growth after an offer is
\[
  \Delta D_{if,t:t+H}^{\mathrm{obs}}
  =
  \Delta D_{i}^{*,H}s_{if}^{0}
  +D_{i}^{*,0}\Delta s_{if}^{H}
  +\Delta D_{i}^{*,H}\Delta s_{if}^{H}
  +u_{iH}.
\]
The first term is growth in total liquidity demand, the second is share
movement, and the third is interaction. In a premium bank, large clients may
grow faster or experience different seasonality than regional-bank clients.
If the model attributes all observed balance growth to $\Delta s_{if}^{H}$, it
will overstate competitor capture. If it attributes all stress-period balance
growth to $D_i^{*}$, it will understate the effect of safety and relationship
share.

The remedy is to estimate the dynamic demand layer and share layer jointly, or
to use a two-step estimator with generated-regressor uncertainty carried
forward. Dynamic-panel methods such as \citet{arellano1991panel} address
persistence and fixed effects, while corporate-liquidity studies motivate why
cash and deposits vary with financing constraints and shocks
\citep{opler1999cash,almeida2004cash,bates2009cash}. Those literatures do not
identify competitor share movement by themselves. Required diagnostics include
forecast error for $D_i^{*}$, share-anchor residuals, and reconciliation of
observed balance changes into move, expansion, and cannibalization. Out-of-time
tests should be stratified by incumbent-bank type, prospect channel, industry,
geography, and relationship size.

\subsection{Campaign Assignment, Relationship-Manager Attention, and Uplift}

The action $o$ is usually not random. Relationship managers contact customers
who look promising, price exceptions are granted to strategically important
firms, and premium banks may selectively pursue clients already likely to move.
Thus
\[
  O_i \not\!\perp\!\!\!\perp
  \{Y_i(o):o\in\mathcal{O}\}\mid \mathcal{I}_{it}
\]
can fail. In that case, propensity, predicted wallet, and treatment effect are
not the same object.

The preferred remedy is randomized or quasi-random learning cells for deposit
offers, pricing, onboarding support, and consolidation campaigns. When random
assignment is not feasible, use overlap diagnostics and orthogonal scores:
\[
\begin{aligned}
  \psi_i(o)
  &=
  m_o(X_i)-m_0(X_i)
  +\frac{1\{O_i=o\}}{e_o(X_i)}\{Y_i-m_o(X_i)\}\\
  &\quad
  -\frac{1\{O_i=0\}}{e_0(X_i)}\{Y_i-m_0(X_i)\}.
\end{aligned}
\]
This score is credible only if overlap holds and the conditioning set contains
the variables that drove assignment. This is the same distinction made in the
uplift and causal-inference literatures between response propensity and
incremental response \citep{rosenbaum1983propensity,lo2002true,wager2018causalforests,chernozhukov2018dml}.
The model should report the effective sample size after weighting, standardized
differences before and after adjustment, placebo effects on pre-treatment
balances, and lift in holdout cells. Otherwise, the model may be used for
prioritization but not causal incrementality.

\subsection{Equilibrium and Competitor Response}

Counterfactual pricing assumes something about competitor behavior. If the
focal bank raises rates for selected SMEs, competitors may match, selectively
retain high-value accounts, or respond through service and credit terms. The
share simulation
\[
  s_{ift}(o)
\]
is therefore conditional on a competitor-response scenario. A unilateral
simulation is useful for small tests but can overstate capture in broad
campaigns.

The remedy is to define counterfactual regimes explicitly:
\[
  s_{ift}^{r}(o,o_{-f}^{0}),\qquad
  s_{ift}^{r}(o,o_{-f}^{\mathrm{match}}),\qquad
  s_{ift}^{r}(o,o_{-f}^{\mathrm{retain}}).
\]
Management should see wallet capture under no-response, rate-match, and
selective-retention scenarios. A recommendation that works only under
no-response assumptions should be treated as a testable hypothesis, not a base
case.

\subsection{Minimum Estimation Standard}

Before the recommended model is used for material targeting or pricing, the
bank should require the following estimation ledger:

\begin{enumerate}[leftmargin=*]
  \item \textbf{Selection ledger:} comparison of current focal-bank customers,
  target prospects, won prospects, and lost prospects; target-population weights
  and sensitivity to unobserved selection.
  \item \textbf{Price ledger:} quote coverage, rejected quotes, exception rates,
  price-imputation accuracy, and evidence on price exogeneity or instruments.
  \item \textbf{Quality ledger:} observed quality proxies, unobserved
  product-market fixed effects, and sensitivity of price elasticities to quality
  controls.
  \item \textbf{Wallet ledger:} source of each competitor-wallet estimate:
  directly observed, anchored, or imputed.
  \item \textbf{Dynamic ledger:} decomposition of balance changes into total
  liquidity demand, share movement, interaction, cannibalization, and residual.
  \item \textbf{Causal ledger:} randomized evidence, quasi-random evidence, or
  associational status for each offer and segment.
  \item \textbf{Counterfactual ledger:} competitor-response scenario, stress
  regime, price assumptions, and uncertainty interval for every policy output.
\end{enumerate}

These ledgers make the estimation problem auditable. They also prevent the
largest practical mistake in this project: presenting a precise-looking wallet
or pricing number when the identifying evidence supports only a ranking, a
sensitivity range, or a learning-cell recommendation.

\section{Validation and Evaluation}

Validation should be separated into ledgers that match the estimands.

\paragraph{Wallet reconstruction.}
Evaluate whether estimated wallet sizes match external validation samples,
customer self-reports, bureau data, account aggregation data, or post-acquisition
observations when a customer consolidates accounts. Metrics include calibration,
rank ordering, interval coverage, and error by industry and firm size.

\paragraph{Product amount and severity.}
Evaluate the conditional amount model separately from product incidence. Metrics
should include balance-weighted error, mean and quantile calibration, prediction
interval coverage, segment-level relative bias, and retransformation diagnostics
for log-dollar models. A high adoption AUC does not validate estimated deposit
balances, loan balances, card spend, or assets under management.

\paragraph{Dynamic deposit demand.}
Evaluate deposit-balance forecasts out of time and across stress regimes.
Metrics should include calibration of prediction intervals, error by industry
and season, stress-period stability, and decomposition of forecast error into
total-liquidity demand error versus focal-bank share error. A model that ranks
large deposit customers well in normal times may still fail if it treats stress
liquidity as a fixed wallet stock.

\paragraph{Stress-condition pricing.}
Evaluate rate sensitivity, deposit beta, and runoff separately by stress regime.
Metrics should include dollar calibration by scenario, rate-elasticity sign and
magnitude stability, branch-share anchor residuals, balance decay after stress,
and NPV sensitivity to funding value and pricing cost. Stress pricing should not
be promoted from normal-period elasticity alone.

\paragraph{Product prediction.}
Evaluate product adoption, balance growth, and sequence prediction using
out-of-time samples. Metrics include AUC, precision at top $k$, calibration,
balance-weighted error, and next-product hit rate. These are predictive metrics,
not proof of causal campaign effectiveness.

\paragraph{Causal response.}
Evaluate campaigns with randomized holdouts or credible quasi-experiments.
Metrics include incremental adoption, incremental balances, risk-adjusted
profit, retention, and adverse outcomes. A campaign should not be promoted based
only on high propensity scores or backtested correlation.

\paragraph{Move-over decomposition.}
Evaluate whether predicted moved amounts match observed evidence about migration
origin in conversion samples. Metrics include cap violations
($\widehat{\Delta a}^{\mathrm{move}}>\widehat{A}^{-f}$), reconciliation error in
the move, expansion, and cannibalization identity; gross and net balance lift;
post-offer decay of migrated balances; and product-level diversion from existing
focal-bank products. Without this validation, the model should report
incremental net balances rather than competitor dollars captured.

\section{Research Gaps}

Several gaps remain for the exact US SME banking use case.

\begin{enumerate}[leftmargin=*]
  \item Most direct cross-selling papers use consumer or general financial
  services data, not modern US SME banking panels.
  \item Competitor product holdings and balances are usually unobserved, so
  validation is difficult.
  \item Even when competitor-held balances are estimated, the movable fraction
  of those balances is usually not identified without conversion, transfer,
  account-aggregation, or experimental evidence.
  \item Deposit size is not merely another product label; it is continuous,
  rate-sensitive, liquidity-sensitive, and operationally tied to payments.
  It should be treated as dynamic demand, not as a constant wallet stock.
  \item The literature on SME-specific deposit dynamics during stress is less
  direct than the broader corporate cash, credit-line, and bank-liquidity
  literature. A bank should validate stress behavior with its own panel data.
  \item Marketing propensity and causal treatment effect are frequently
  conflated in practice.
  \item Relationship lending creates joint value across deposits, payments,
  lending, and monitoring; product-by-product models may understate this
  complementarity.
  \item Product-assortment, next-best-product, and bundle design are useful
  extensions, but they are not the main problem addressed here. They should be
  treated as further research after the deposit-wallet and deposit-capture
  estimands are validated.
  \item Compliance, fair-lending, privacy, and model-risk constraints are
  central in bank deployment but are not the main subject of the academic
  cross-sell literature.
\end{enumerate}

\section{Recommended Implementation Roadmap}

A bank can proceed in stages.

\paragraph{Stage 1: Define the wallet taxonomy.}
Start with a deposit-centered taxonomy: operating deposits, savings and
liquidity products, treasury-linked balances, and externally held competitor
balances. Define incidence, size, competitor-held amount, movable amount, net
incremental amount, runoff, and profitability measures. Other products can be
added as relationship context or later extensions.

\paragraph{Stage 2: Build baseline wallet and liquidity proxies.}
Use internal balances, transaction flows, industry benchmarks, revenue
estimates, and market-level deposit data to create transparent baseline wallet
estimates. For deposits, construct dynamic liquidity proxies such as trailing
cash inflow volatility, payroll coverage months, working-capital intensity,
seasonality, and local market position. These baselines are necessary
comparators for later machine-learning models.

\paragraph{Stage 3: Fit deposit incidence and size models.}
Use two-part models for deposit relationships: an incidence model for whether
the SME has relevant deposits and a severity model for positive balances. Use
other products as covariates or latent relationship context. Report uncertainty
intervals, conditional-amount calibration, quantile calibration, and
segment-level dollar bias.

\paragraph{Stage 4: Estimate dynamic deposit demand.}
Fit panel time-series models for deposit balances and total liquidity demand.
Include lagged deposits, transaction flows, industry-by-time effects,
geography-by-time effects, rate environment, local competition, and relationship
terms. Benchmark against seasonal naive and fixed-effects baselines before using
more complex sequence models.

\paragraph{Stage 5: Estimate deposit pricing and share response.}
Use rate, fee, service, branch-market, and relationship variation to estimate
deposit share response. Treat branch-level market share as a noisy aggregate
anchor when it cannot separate consumer, SME, commercial, and wholesale
deposits.

\paragraph{Stage 6: Stress-test deposit demand.}
Estimate stress scenarios for balance growth, cash burn, flight-to-safety
migration, and post-stress runoff. Validate these scenarios with the bank's own
panel history and management judgment.

\paragraph{Stage 7: Introduce randomized campaign holdouts.}
For each major deposit action, maintain holdouts and estimate incremental lift.
Avoid using raw purchase propensity as the promotion criterion.

\paragraph{Stage 8: Estimate move-over and cannibalization.}
Use conversion events, account-transfer evidence, payoff and balance-transfer
records, account aggregation, direct-deposit or ACH migration, and randomized
consolidation offers to estimate the fraction of competitor-held balances that
move under each offer. Separate migrated competitor wallet from total-demand
expansion and cannibalization of existing focal-bank balances.

\paragraph{Stage 9: Optimize for risk-adjusted relationship value.}
Combine expected incremental balances, fee revenue, interest margin, funding
value, credit risk, servicing cost, liquidity value, and compliance constraints.

\paragraph{Stage 10: Monitor drift and fairness.}
Monitor calibration, response, customer complaints, adverse action patterns,
industry/geographic stability, stress-regime performance, and
relationship-manager override behavior.

\section{Conclusion}

There is substantial literature relevant to SME deposit wallet-size estimation,
but it is distributed across marketing science, banking economics, SME finance,
and causal inference. The direct cross-selling literature provides models for
latent product need and wallet context, while share-of-wallet research clarifies
the distinction between customer retention and relationship depth.
Deposit-demand and bank-competition research explain why deposits require
special treatment. Corporate liquidity and credit-line research shows why
deposits should not be treated as a fixed-size wallet bucket: firms adjust cash
balances and liquidity demand as conditions change. Relationship-lending papers
show why SME transaction accounts and credit products are jointly valuable.
Causal uplift methods are needed to decide which deposit actions will actually
move wallet.

The recommended solution is therefore not a single score. It is a modular
decision system: latent wallet estimation, dynamic deposit demand,
focal-bank share response, deposit pricing, campaign uplift, and constrained
profit optimization. Such a system should be validated with external or
post-conversion wallet evidence, branch-market-share anchors used cautiously,
out-of-time deposit forecasts, stress-regime diagnostics, and randomized
marketing holdouts before it is used for major pricing or targeting decisions.

\section*{Scope Note}

This document is a technical literature synthesis for problem formulation and
model design, not a full systematic review. It emphasizes sources that are
directly useful for banking cross-selling, wallet share, deposit competition,
SME relationship lending, dynamic liquidity demand, and causal response. Before
using the document as a formal literature-review artifact, the bank should
complete a documented source-support ledger, backward and forward snowballing,
retraction/erratum checks, and an omission-risk register.

\bibliographystyle{plainnat}
\bibliography{sme_wallet_refs}

\end{document}
